\documentclass{article}
\usepackage{jheppub}
\usepackage{graphicx} % Required for inserting images
\usepackage{hyperref}
\usepackage{subcaption}
\usepackage{amsmath, amssymb}
\usepackage{bbm}
\usepackage{xcolor}
\usepackage{tikz}
\usepackage{tikz-3dplot}
\usepackage[normalem]{ulem}
\usepackage{braket}
\usepackage{booktabs,tabularx}
\newcommand{\Tr}{\mathrm{tr}}
\newcommand{\KW}{\mathrm{KW}}
\newcommand{\mA}{\mathcal{A}}
\newcommand{\mB}{\mathcal{B}}
\newcommand{\mC}{\mathcal{C}}
\newcommand{\mL}{\mathcal{L}}
\newcommand{\mH}{\mathcal{H}}
\newcommand{\mT}{\mathcal{T}}
\newcommand{\nn}{\nonumber}
\newcommand{\lb}{\left(}
\newcommand{\rb}{\right)}
\newcommand{\ep}[]{\epsilon}
\usetikzlibrary{patterns.meta}
\usetikzlibrary{fadings}
\usepackage{graphicx}
\usepackage{amsmath}
\usepackage{amssymb}
\usepackage{hyperref}
%\usepackage{showlabels}
\usepackage{amsthm}
\newtheorem{theorem}{Theorem}
\newtheorem{lemma}[theorem]{Lemma}
\newtheorem{corollary}[theorem]{Corollary}
\newtheorem{proposition}[theorem]{Proposition}
\theoremstyle{definition}
\newtheorem{definition}{Definition}[section]
\newtheorem{prop}{Conjecture}
\newtheorem{cor}{Corollary}
%-------------------packages-------------------------
\usepackage{enumerate}
\usepackage{tikz}

\bibliographystyle{ieeetr}

\definecolor{sulaiman}{rgb}{255,0,255}
\newcommand{\SA}[1]{{\color{sulaiman}\footnotesize{(SA) #1}}}
\newcommand{\NL}[1]{\textcolor{blue}{#1}}

\title{Non-invertible symmetry is non-unitarily correctable QEC}
\author{ }
\date{}

\abstract{We argue that known examples of non-invertible symmetries can be recast in terms of quantum error correcting codes.  Specifically, we view a given non-invertible symmetry as an error map acting on an algebra of observable operators.  For examples of internal symmetries that do not mix with spacetime transformations, we can always define a conditional expectation that maps to the passively correctable subalgebra. For operators charged under the non-invertible symmetry, we find that the recovery map is given by the non-unitary dual superoperator.}

\begin{document}

\maketitle
Key:

\color{blue} blue \color{black} = calculations or details

\color{brown} brown \color{black} = will add

\SA{purple} = comments/confusions/questions
\section{Research questions}

\begin{itemize}
  \item Is any non-invertible symmetry a non-unitarily correctable Quantum Error Correction Code (QECC)? Is a local QEC code? If yes, can one apply the tools of QECC to shed light on non-invertible symmetry?
  \item Can every non-unitarily correctable {\it local} QECC be formulated as a non-invertible symmetry? If yes, can one apply recent developments in non-invertible symmetry to shed light on QECC?
  \item In QEC, it is known that the use of operator algebraic language leads to a simplification and provides a generalization route to the continuum limit; i.e., algebraic quantum field theory (AQFT). Can one use operator algebraic tools of AQFT to shed light on non-invertible symmetries in QFT?
\end{itemize}
To start: Start with the simplest example of a non-invertible symmetry and show that it is a non-unitarily correctable QECC.

Here are the simplest examples of non-invertible symmetries:
\begin{itemize}
  \item {\bf Discrete:} KW model, Rep $D_8$
  \item {\bf Contiuum:} 1+1D Ising CFT, massless chiral fermions in QED$_4$ (Chiral anomaly as non-invertible symmetry).
\end{itemize}

% \subsection{To do list}

% \begin{enumerate}
%     \item\textit{Rep $D_8$}
%     \subitem Work through details of Rep $D_8$, show that the ''charge" of the noninvertible operator cannot add.

%     \item\textit{Continuum Limit of TFIM, Ising CFT}
%     \subitem Understand if the operator algebra generated by the continuum limit of 1D lattice falls under the main conclusion of the paper, or if the machinery of CFT is needed
% \end{enumerate}

\color{brown} In introduction or conclusion should add comment about destabilizer vs. stabilizer, i.e. could have chosen the framework where the NIS protects the qubits instead of acts as noise.\color{black}

\section{Inclusion of Algebras as Generalization of groups}

Mathematically speaking, one framework is to start with an inclusion of algebras with finite index $\mA\subset \mB$. This data is enough to define a finite semi-simple Hopf algebra and a Q-system (unitary Frobeinus algebra). Longo shows that this data is dual to the category of endomorphisms of $\mA$ \cite{longo1994duality}. If we take the endomorphisms of $\mA$ to be localizable and transportable we obtain the DHR theory.

\section{Crossed product $\mA\rtimes_\alpha G$}

Consider the action of a group $\alpha:G\to \text{Aut}(\mB)$. Of course, for each $g\in G$ the group action $\alpha_g$ is an automophism of the algebra. Consider the trivial automorphism $\alpha_1(a)=a$ and $\alpha_g(a)$. We define an intertwiner between the two:
\begin{eqnarray}\label{intertwineG}
  R_g \alpha_1(a)=R_g a=\alpha_g(a) R_g\ .
\end{eqnarray}
Applying these intertwiners successively gives
\begin{eqnarray}
  R_{g_1}R_{g_2}a=(\alpha_{g_1}\circ\alpha_{g_2})(a)R_{g_1}R_{g_2}=\alpha_{g_1g_2}(a)R_{g_1}R_{g_2}
\end{eqnarray}
which implies that we can identity
\begin{eqnarray}
  R_{g_1g_2}=R_{g_1}R_{g_2}\ .
\end{eqnarray}
In other words, the composition of endormophisms, implies a multiplication rule for the intertwiners. In the case of a group $\alpha_g(a)=U_g a U_g^\dagger$, the intertwiners is the simply the unitary $R_g=U_g$. $R_g$ has the interpretation of a defect because as you pass the operators through it, they transform by the group action.

An element of the crossed product can be viewed as the pair $(a,R_g)$ or simply $a R_g$. Then, the intertwiner equation (\ref{intertwineG}) says that there is a natural product for the elements of the crossed product
\begin{eqnarray}
  a_1 R_{g} a_2 R_h= a_1 \alpha_g(a_2) R_g R_h= a_1\alpha_g(a_2) R_{gh}\ .
\end{eqnarray}
The crossed product elements $a U_g$ represent the twisted sector $g$.

The simplest example of the $\mathbb{Z}_2$ symmetry of 2d Ising CFT. The algebra of operators is generated by $\{1,\epsilon,\sigma\}$ Virasoro primaries. The symmetry automorphisms are $\{\alpha_+=\text{id},\alpha_-\}$ with $\alpha_-(1)=1$, $\alpha_-(\epsilon)=\epsilon$ and $\alpha_-(\sigma)=-\sigma$. The symmetry is generated by the charge operator $\mathcal{L}_\epsilon$ and it satisfies the analog of intertwiner equations (\ref{intertwineG}):
\begin{eqnarray}
  \mL_\epsilon 1=1\mL_\ep, \qquad \mL_\ep \epsilon=\epsilon \mL_\ep, \qquad \mL_\ep\sigma=-\sigma \mL_\ep,\qquad \mL_\ep \mL_\ep=1\ .
\end{eqnarray}
This immediately implies that the twisted sector operators are
\begin{eqnarray}
  \alpha_-(a)=\mL_\ep a \mL_\ep\ .
\end{eqnarray}

Note that the group action is an outer automorphism because $\mu\notin \mA$, and that is why we are calling it a twist and not a local operator.

The crossed product space $\mA\rtimes_\alpha G$ includes both the vacuum sector $a 1$ and the twisted sector as $a$ and $a \mu$. An alternative point of view is to take the direct sum of the untwisted operators $\alpha_+(\mA)=\mA$ with the twisted ones $\alpha_-(\mA)=\mu \mA \mu$ which can be viewed as the map $\mA\to \mA^{\otimes |G|}$:
\begin{eqnarray}
  \lambda=\mA\otimes G=\oplus_{g\in G}\alpha_g(\mA)\ .
\end{eqnarray}
More precisely, $\lambda=\mA\otimes l^2(G)$ is an object similar to the regular representation of $G$.
\NL{Is the folowing correct? For an infinite dimensional algebra $\mA$ and a finite group $G$, the map $\mA\to \mA\otimes G$ can be an isomorphism.}

The invariant algebra of $\mA$ is found by taking the group average (conditional expectation)
\begin{eqnarray}
  &&\mathcal{E}_G(a)=\frac{1}{|G|}\sum_{g\in G}\alpha_g(a)=\frac{1}{2}(a+\mu a \mu)\nn\\
  &&\mA^G=\mathcal{E}_G(\mA)
\end{eqnarray}
This projects the algebra $\mA$ down to the subalgebra $\mA^G$ of $\{1,\epsilon\}$ projecting $\sigma$ out. We end up for two inclusions 
\begin{eqnarray}
  \mA^G\subset \mA\subset \mA\rtimes_\alpha G\ .
\end{eqnarray} 
Each $\alpha_g$ is an embedding of $\mA$ inside the larger algebra $\lambda$. 

Restricting the map $\lambda$ above to the invariant algebra $\mA^G$ gives
\begin{eqnarray}
  \oplus_{g\in G} \mu \mA^G \mu=\oplus_{g\in G}\mA^G
\end{eqnarray}
Of course, if we push invariant operators through $\mu$, we are going to find $\mu a \mu=a$, however, more abstractly we can view each $R_g^\dagger a R_g$ as a representation of $\mA^G$, and more generally for a non-Abelian group the restriction of $\lambda$ to $\mA^G$ gives
\begin{eqnarray}
  \lambda\big|_{\mA^G}=\oplus_{r \in \hat{G}}d_G(r)\pi_r\ .
\end{eqnarray}

The map $\lambda$ can be viewed as the canonial endomorphism $J_{\mA^G}J_\mA$ of the inclusion $\mA^G\subset \mA$ 

\subsection{Dual group}

Elements of the dual group $\chi\in \hat{G}$ are the irreducible representations of $G$. In the case of Abelian group, the dual goup is itself a group, and has a natural action on the crossed product $\mA\rtimes_\alpha G$:
\begin{align}
  \hat{\alpha}:\hat{G}\to \text{Aut}(\mA\rtimes_\alpha G), \qquad \hat{\alpha}_\chi(a R_g) = \chi(g) a R_g
\end{align}
For example, in the case that $G=\mathbb{Z}_n$ labled by $g=1,\cdots, n$, the dual group is labled by $\hat{G}=\mathbb{Z}_n$ with $\hat{g}=1,\cdots, n$, and the dual action is 
\begin{align}
&\hat{\alpha}_k(a R_g) = e^{2\pi i k g/n} a R_g  \nn\\
&V_k a R_g=e^{2\pi ik g/n} a R_gV_k\ .
\end{align}
 Intuitively, we can view $V_k$ as commuting with $a$ and satisfying $V_k R_g=e^{2\pi i k g/n} R_g V_k$. 
 
Note that in this case, the dual group is a group $V_{\hat{g}_1\hat{g}_2}=V_{\hat{g}_1}V_{\hat{g}_2}$, and the dual action is an inner automorphism of the crossed product algebra.



\section{Invertible symmetry and group action on algebras}

The observables of a quantum system form a $C^*$-algebra, and invertible symmetry transformations are automorphisms of this algebra \footnote{An automorphism of a $C^*$ algebra is $*$-isomorphism from the algebra to itself.}. Since automorphisms form a group, to study invertible symmetries in quantum mechanics, one needs to study the action of a group on a $C^*$ algebra which is captured by the notion of a {\it reversible $C^*$-dynamical system}.

More generally, non-invertible and generalized symmetries correspond to endomorphisms of the observable algebra ($*$-homomorphisms from the algebra to itself). Since endomorphisms form a monoid (semigroup with identity element), to study non-invertible symmetries, we have to generalize beyond groups and discuss {\it irreversible dynamical system}. Here, we start with a review of a reversible dynamical system and its connection to invertible symmetry, and postpone the discussion of non-invertible symmetries to the next Section.

\subsection*{Reversible $C^*$-dynamical system, crossed product}

A reversible $C^*$-dynamical system is a trip $(\mA, G, \alpha)$ where $\mA$ is a $C^*$-algebra, $G$ is a locally compact group, and $\alpha:G\to \text{Aut}(\mA)$ is a group homomorphism:
\begin{eqnarray}
  \alpha_{g_1}\circ\alpha_{g_2}(a)=\alpha_{g_1g_2}(a)\ .
\end{eqnarray}
We choose $G$ to be locally compact so that we have a left and right invariant Haar measure. To be able to integrate elements of $\mA$ with respect to this measure, we often require the action of $\alpha$ to be strongly continuous in $g$:
\begin{eqnarray}
  \lim_{g\to e}\|\alpha_g(a)-a\|=0\ .
\end{eqnarray}
where $e$ is the identity element of the group, and the continuity in the $C^*$-norm topology.

In physics, we are interested in covariant representations of a $C^*$-dynamical system. A covariant representation of a $C^*$-dynamical system $(\mA,G,\alpha)$ is a trip $(\pi, U,\mH)$ where $\pi:\mA \to B(\mathcal{H})$ is a non-degenerate\footnote{A representation is non-degenerate if there exists no $\ket{\psi}\in\mathcal{H}$ such that $a\ket{\psi}=0$ for all $a\in\mA$.} $*$-representation of the algebra, and $U:G\to U(\mathcal{H})$ is a unitary representation of the group:
\begin{eqnarray}
  \forall g\in G, a\in\mathcal{A}:\qquad \pi_g(a)=U_g \pi(a)U_g^\dagger\ .
\end{eqnarray}
\begin{theorem}
  There is a one-to-one correspondence between covariant representations $(\pi,U,\mathcal{H})$ of a $C^*$-dynamical system, and non-degenerate $*$-representations $\tilde{\pi}$ of the crossed product $\mA\rtimes_\alpha G$.
\end{theorem}
Given a representation $\mathcal{H}_0$ of a $C^*$-algebra, we can construct a representation of the crossed product $L^2(G,\mH_0)$ and obtain a covariant representation. However, this depends explicitly on $\mH_0$, hence is non-uniquely, and often this Hilbert space is too large for physics.

\subsection*{Reversible von Neumann dynamical system}

\begin{definition}[W*-Dynamical System]
  A $W^*$-dynamical system is a triple $(M, G, \alpha)$ where:
  \begin{enumerate}
    \item $M$ is a von Neumann algebra acting on a Hilbert space.
    \item $G$ is a locally compact group.
    \item $\alpha: G \to \mathrm{Aut}(M)$ is a group homomorphism such that for every $x \in M$ and $\phi \in M_*$, the map $g \mapsto \phi(\alpha_g(x))$ is continuous (point-$\sigma$-weak continuity).
  \end{enumerate}
\end{definition}
We say we have an {\it invertible symmetry} with group $G$ when we have a $W^*$-dynamical system.

\begin{theorem}
  Let $(M, G, \alpha)$ be a W*-dynamical system and let $(M, \mathcal{H}, J, \mathcal{P})$ denote the unique standard form of $M$. There exists a unique strongly continuous unitary representation $g \mapsto U_g$ of $G$ on $\mathcal{H}$ such that for all $g \in G$:
  \begin{align}
    U_g x U_g^* &= \alpha_g(x) \quad \forall x \in M, \\
    J U_g &= U_g J, \\
    U_g \mathcal{P} &= \mathcal{P}.
  \end{align}
\end{theorem}

\begin{proof}
  The proof relies on the uniqueness of the standard form of a von Neumann algebra (Haagerup). See \cite{haag2012local}. Fix $g \in G$. Since $\alpha_g$ is a *-automorphism of $M$, we may view it as an isomorphism between two copies of the standard form: $(M, \mathcal{H}, J, \mathcal{P}) \xrightarrow{\alpha_g} (M, \mathcal{H}, J, \mathcal{P})$.
  By the uniqueness theorem of the standard form, for any automorphism $\beta$ of $M$, there exists a unique unitary $V$ on $\mathcal{H}$ implementing $\beta$ (i.e., $V x V^* = \beta(x)$) such that $V J V^* = J$ and $V \mathcal{P} = \mathcal{P}$. Applying this to $\beta = \alpha_g$, there exists a unique unitary $U_g$ satisfying conditions (1), (2), and (3).

  \paragraph{Group Property:}
  Consider elements $g, h \in G$. The product $U_g U_h$ is unitary and satisfies:
  \begin{enumerate}
    \item $(U_g U_h) x (U_g U_h)^* = U_g (\alpha_h(x)) U_g^* = \alpha_g(\alpha_h(x)) = \alpha_{gh}(x)$,
    \item $U_g U_h J = U_g J U_h = J U_g U_h$,
    \item $U_g U_h \mathcal{P} = U_g \mathcal{P} = \mathcal{P}$.
  \end{enumerate}
  By the uniqueness established above, $U_{gh}$ is the \textit{only} operator satisfying these conditions for the automorphism $\alpha_{gh}$. Therefore, $U_{gh} = U_g U_h$.

  To show strong continuity, it suffices to show $\lim_{g \to e} \|U_g \xi - \xi\| = 0$ for $\xi \in \mathcal{P}$. For any $\xi \in \mathcal{P}$, the map $x \mapsto \langle \xi, x \xi \rangle$ defines a normal positive functional $\omega_\xi \in M_*$. Standard estimates on the natural cone (specifically the Powers-St{\o}rmer inequality generalized by Araki) bound the vector norm distance by the predual norm distance:
  $$ \|U_g \xi - \xi\|^2 \leq \| \omega_\xi \circ \alpha_g - \omega_\xi \|_{M_*}. $$
  Since $\alpha$ acts continuously on the predual $M_*$, the right-hand side vanishes as $g \to e$. Thus $U$ is strongly continuous.
\end{proof}

Reversible symmetry transformations act as physical unitary transformations $U_g$ within the same Hilbert space.

\section{Endomorphisms of algebras and non-invertible symmetries}

If we generalize by taking the range of $\alpha:G\to \text{End}(\mA)$ to be the endomorphisms ($*$-homomorphisms from $\mA$ to $\mA$) we obtain a {\it non-invertible dynamical system}.

\subsection*{Representation of $C^*$-algebra}

A representation of an $C^*$-algebra is a $*$-homorphism $\pi:\mathcal{A}\to \mathrm{End}(\mathcal{H})$. Physicists are familiar with the statement that for a group $G$, given a pair of representations: $\pi_1(U)$ and $\pi_2(U)$ for $U\in G$, the tensor product $\pi_1(U)\otimes \pi_2(U)$ and the direct sum of representations $\pi_1(U)\oplus \pi_2(U)$ are also representations. For algebras, the direct $\pi_1(\mathcal{A})\oplus \pi_2(\mathcal{A})$ is a representation; however, in general $\pi_1(\mathcal{A})\otimes \pi_2(\mathcal{A})$ is not a representation! The reason is that in addition to multiplication, an algebra also has an addition operation, and the naive map $\pi(a)\otimes \pi(a)$ violates linearity. If the algebra has a co-product operation: $\Delta:\mathcal{A}\to \mathcal{A}\otimes \mathcal{A}$ (it is a bi-algebra), then we can fuse two representations to obtain a new representation using the co-product:
\begin{eqnarray}
  \pi_{\mathcal{H}_1\otimes \mathcal{H}_2}(a)=(\pi_{\mathcal{H}_1}\otimes \pi_{\mathcal{H}_2})\Delta(a)\ .
\end{eqnarray}
A key example of bialgebras in physics is the universal enveloping algebra of a Lie algebra.

\subsection*{Endormophisms of $C^*$-algebra}

We are interested in the endomorphisms ($*$-homorphisms) of the local algebra of observables of a general quantum system:
\begin{eqnarray}
  &&\rho:\mathcal{A}\to \mathcal{A}\nn\\
  &&\rho(a_1)\rho(a_2)=\rho(a_1a_2)\nn\\
  &&\rho(a^*)=(\rho(a))^*\ .
\end{eqnarray}
A representation is $*$-homorphism $\pi:\mA\to \mathrm{End}(\mathcal{H})$. Therefore, an endomorphism $\rho$ can be viewed as a way of ``twisting" a representation $\pi_0$ into a new one $\pi_\rho$:
\begin{eqnarray}
  \pi_\rho(a)=\pi_0(\rho(a))\ .
\end{eqnarray}
In other words, $\rho$ acts as a {\it functor} on the category of representations of $\mA$. It takes a representation $(\pi,\mathcal{H})$ to a new representation $(\pi\circ\rho,\mathcal{H})$.
Since we want endomorphisms $\rho$ to be a functor, this means that the direct sum of representations should induce an addition rule for endomorphisms: $\rho_1\oplus \rho_2$:
\begin{eqnarray}
  \pi_{\rho_1\oplus \rho_2}(a)=\pi_1(a)\oplus\pi_2(a)\ .
\end{eqnarray}
To arrange for this direct-sum we need linear isometries $V_1^\dagger V_1=V_2^\dagger V_2=1$ such that $V_1V_1^\dagger+V_2V_2^\dagger=1$:
\begin{eqnarray}
  (\rho_1\oplus\rho_2)(a)=V_1\rho_1(a)V_1^\dagger+V_2\rho(a)V_2^\dagger\ .
\end{eqnarray}
The condition that $V_i$ are isometries follows from the fact that $(\rho_1\oplus\rho_2)$ is a homomorphism, i.e. sends multiplication to multiplication. The condition that $V_1V_1^\dagger+V_2V_2^\dagger=1$ follows from the requirement that $(\pi_1\oplus\pi_2)(1)=1$. The elements $V_i$ generate a {\it cuntz} algebra. Note that there are no realizations of $V_i$ isometries in finite-dimensional matrix algebras.

We call an endomorphism $\rho$ simple if it cannot be decomposed into a direct sum of other endomorphisms. If $\pi_0$ is a faithful irreducible representation and $\rho$ is a simple endomorphism, then $\pi\circ \rho$ is also an irreducible representation. We use the notation $\rho_r$ and $\pi_r$ to denote simple endomorphisms and their corresponding representations.

Consider the composition of two simple endomorphisms $\rho_{r_1}\circ \rho_{r_2}$ is not simple. Often, we can decompose the corresponding representation into irreducible ones:
\begin{eqnarray}
  \pi(\rho_{r_1}\circ\rho_{r_2}(a))= \oplus_r (\rho_r(a)\otimes 1_{N^r_{r_1r_2}})
\end{eqnarray}
where $N^r_{r_1,r_2}$ are called the multiplicities of representation $r$. This induces the relations
\begin{eqnarray}
  &&(\rho_{r_1}\circ \rho_{r_2})(a)=\sum_k \sum_{\alpha=1}^{N^k_{r_1,r_2}}W_{k,\alpha} \rho_k(a)W_{k,\alpha}^*\nn\\
  &&\rho_{r_1}\times \rho_{r_2}=\sum_r N^r_{r_1r_2} \rho_r\ .
\end{eqnarray}

We call the dual of an endomorphism a {\it transfer operator}.
\begin{definition}[Transfer Operator]
  If $\alpha: \mathcal{A} \to \mathcal{A}$ is an endomorphism, a transfer operator $\mathcal{L}: \mathcal{A} \to \mathcal{A}$ is a positive linear map such that:
  \begin{equation}
    \mathcal{L}(\alpha(a)b) = a \mathcal{L}(b) \quad \forall a,b \in \mathcal{A}
  \end{equation}
\end{definition}
Setting $b=1$, this implies $\mathcal{L}(\alpha(a)) = a$. If $\alpha$ represents a deterministic evolution that loses information, $\mathcal{L}$ represents an ``average'' over the pre-images. It tracks how densities (states) evolve backward against the flow of observables.

\subsection*{Non-invertible symmetry}

To formalize a non-invertible symmetry, one can define an $C^*$ semi-dynamical system that is irreversible as
\begin{definition}[C*-Semi-Dynamical System]
  A \textbf{C*-Semi-Dynamical System} is a triple $(\mathcal{A}, S, \alpha)$ where:
  \begin{enumerate}
    \item $\mathcal{A}$ is a C*-algebra.
    \item $S$ is a semigroup with unit $e$ (monoid) (e.g., $\mathbb{N}$ or $\mathbb{R}_+$).
    \item $\alpha: S \to \mathrm{End}(\mathcal{A})$ is a homomorphism:
      \begin{equation}
        \alpha_{g_1} \circ \alpha_{g_2} = \alpha_{g_1g_2}, \quad \alpha_e = \mathrm{id}
      \end{equation}
  \end{enumerate}
\end{definition}
More generally, we can generalize beyond semi-groups and consider any set $S$. We say we have a {\it non-invertible symmetry} if we have a map $\alpha:S\to \mathrm{End}(\mathcal{A})$. Of course, we need a notion of multiplication to make sense of successive applications of the transformation, and consider homomorphisms
\begin{eqnarray}
  \alpha_{g_1}\circ \alpha_{g_2}=\alpha_{g_1 g_2}, \alpha_e=\text{id}
\end{eqnarray}
with an identity element. Assuming that multiplication is associative makes $S$ a monoid, and we are back to the $C^*$-semi-dynamical system. To go beyond semi-groups, we replace $S$ with an algebra, so that we have multiplication and addition. For now, we will assume that multiplication is commutative, meaning
\begin{eqnarray}\label{muliplication}
  \forall \rho_1,\rho_2\in S:\qquad \rho_1\circ \rho_2=\sum_i c_i \rho_i
\end{eqnarray}
where $c_i$ take values in a field (say, natural numbers $\mathbb{N}$, or complex numbers $\mathbb{C}$).

In the GNS Hilbert space, we pick the representation $\pi_0(a)$. Then, for any endomorphism $\rho$ we find a new representation $\pi_\rho(a)$, and {\it charged intertwiners} $V_\rho$ that intertwine the two representations:
\begin{eqnarray}
  V_\rho \rho_0(a)=\pi_\rho(a) V_\rho\ .
\end{eqnarray}
Often, a representation $\rho$ can be decomposed into irreducibles
\begin{eqnarray}\label{decompose}
  \pi_\rho(a)=\oplus_r \oplus_{i=1}^{d_r(\rho)}\pi_r(a)=\oplus_r (\pi_r(a)\otimes 1_{d_r})
\end{eqnarray}
where $d_r$ are positive integers called multiplicities. Therefore, elements of $S$ can be labeled by representation labels $\rho$, with $e$ corresponding to the initial vacuum GNS representation $\pi_0$:
\begin{eqnarray}
  \alpha_\rho(a)\equiv \rho(a)\ .
\end{eqnarray}

Comparing (\ref{muliplication}) with (\ref{decompose}) makes it clear that the multiplication in the set $S$ is the tensor product of representations:
\begin{eqnarray}
  &&    \pi((\alpha_{r_1}\circ\alpha_{r_2})(a))=\oplus_r (\pi_r(a)\otimes 1_{N^r_{r_1,r_2}})
\end{eqnarray}
where

A major theme in operator algebras is \textbf{Dilation Theory}. If you have an irreversible system $(\mathcal{A}, \alpha)$, can you find a larger algebra $\mathcal{B}$ and an automorphism $\beta$ on it, such that $\mathcal{A}$ lives inside $\mathcal{B}$ and $\alpha$ is just the restriction of $\beta$?

\begin{theorem}[Minimal Dilation]
  Under mild conditions, any C*-dynamical system with endomorphisms can be embedded into a larger C*-dynamical system with automorphisms:
  \begin{equation}
    \mathcal{A} \subset \mathcal{B}, \quad \beta_t|_{\mathcal{A}} = \alpha_t \quad (\text{for } t \ge 0)
  \end{equation}
\end{theorem}

This is the operator-algebraic version of ``Purifying the system'' or finding the ``Naimark Dilation.'' It suggests that irreversibility is often an artifact of looking at a subsystem.

{\bf Cuntz algebra:}

The prototype for this structure is not the rotation algebra (torus), but the \textbf{Cuntz Algebra} $\mathcal{O}_n$.

Consider a Hilbert space with an infinite basis. You can define ``shift'' operators (isometries) $S_i$ that are essentially endomorphisms of the space. Let $S_1, \dots, S_n$ be isometries ($S_i^* S_i = 1$) such that their ranges partition the space ($\sum S_i S_i^* = 1$).

We can define an endomorphism $\Phi$ on operators $X$ by:
\begin{equation}
  \Phi(X) = \sum_{i=1}^n S_i X S_i^*
\end{equation}

Instead of unitary generators $U_g$ implementing the action, we get \textbf{Isometry} generators $V_s$ such that:
\begin{equation}
  V_s a V_s^* = \alpha_s(a)
\end{equation}
Since $V_s$ is an isometry ($V_s^* V_s = 1$) but not a unitary ($V_s V_s^* < 1$), this reflects the irreversible nature of the system.

\section{Categorification of endomorphisms}

Use this reference to review \cite{bischoff2015tensor}

\section{Fusion $n$-categories, endomorphisms of local algebras and symmetries}

\section{Introduction to non-invertible symmetry}

\SA{Given the new structure of the paper, I don't know if we need section 8, maybe the content of section 8.2 can be moved to an introduction but happy with whatever you think is best.}


\subsection{Motivating examples}


In this section, we start with two simple 1D examples of non-invertible symmetry: 1) Kramers-Wannier (KW) in 1D lattice 2) $D_8$-lattice.
\color{brown}
List the key features and shortcomings of each example

Shu-Heng: The advantage of the first is that once you generalize to higher dimensions, it results in a $0$-form (global) non-invertible symmetry, whereas the higher-dimensional generalization of $D_8$ results in a higher-form non-invertible symmetry.

\color{black}


\subsection{General theory of non-invertible symmetry}

Non-invertible symmetries can be understood using the unified framework of generalized symmetries~\cite{gaiotto2015generalizedglobalsymmetries} \color{brown} add citations of Shao, Schafer-Nameki, etc. reviews \color{black}.  As seen above, some of the simplest examples of non-invertible symmetries appear in lattice Hamiltonian systems.  While the theory of generalized global symmetries was originally conceived in the context of relativistic quantum field theory, we can still apply central principles to the lattice \cite{seiberg2024noninvsymmetriestensorproductspace}.  Specifically, to every symmetry operator we associate a corresponding defect.  We can think about symmetry operators as being localized in time and defect operators being localized in space.  In contrast with relativistic theories, lattice Hamiltonian operators and defects are not interchangeable.  One way to represent the action of defects in this setting is to modify the Hamiltonian locally.  This will be the most natural viewpoint for our purposes.

Wigner's theorem is derived with the assumption that symmetries preserve transition amplitudes.  If we relax this condition but still require locality and commuting with the Hamiltonian, we can expand our classification of symmetries to include non-invertible operators such as $\textrm{KW}$ and $D$ above.  These operators have many of the useful properties we would expect of ordinary symmetries \color{brown} add citations for phase diagram, rg flow constraints, selection rules \color{black}.  Furthermore, they nicely fit into the operator-defect framework described above. \color{brown}add citation for generalized Wigner's theorem \color{black}

\section {Non-invertible symmetry as non-unitarily correctable QECC}
\label{sec:IntrotoNUCNIS}


% Algebraically, it is more natural to think of the subalgebra $\tilde{\mA}$ as the algebra generated by $X_j$ with $j\in S$ and identity, and the subalgebra $\tilde{\mB}$ as the algebra generated by $Z_jZ_{j+1}$ with $j\in S$ and identity.

% \subsection{Operator algebraic QEC}

% \SA{Maybe better to format here as: Symmetries and QEC as follows:}

\subsection{Symmetries and QEC}

We now introduce the connection between symmetries of an operator algebra and exact quantum error correction codes.  Throughout this section, we will refer to the example of Krammers Wannier duality on the lattice to illustrate the salient features of non-invertible symmetries as noise maps.  The readers are encouraged to review Appendix~\ref{sec:LatticeKW} for an algebraic approach to the example.

\subsubsection{Invertible Global Symmetry as exact QECC}

Let us now discuss how a an ordinary 0-form symmetry can be viewed as a model of quantum error correction.  First, it is helpful to review some notions of algebraic QEC, details of which can be found in Ref.~\cite{furuya2022real}.  In the Heisenberg picture of quantum error correction, a quantum error correcting code (QECC) is given by the tuple $(\iota, \mathcal{R}, \Phi, \alpha)$ along with operator algebras $\mathcal{B}$ and $\mathcal{A}$. Here, $\iota:\mB\to \mA$ is an isometric encoding map that takes operators in our code algebra $\mB$ to the algebra of physical operators $\mA$. $\alpha:\mA\to \mB$ is our decoding map. The recovery map is $\mathcal{R}:\mA\to \mA$ and the error map $\Phi:\mA\to \mA$ is a unital and completely postive (CP). As CP maps, they have the Kraus representations
\begin{eqnarray}\label{Krausforms}
&&\alpha(a)=W^\dagger a W,\qquad \iota(b)=WbW^\dagger\\\
&&\Phi(a)=\sum_r V_r^\dagger a V_r,\qquad \mathcal{R}(a)=\sum_r R_r^\dagger a R_r\ .
\end{eqnarray}

We have an error correction if for all the code operators $b\in \mB$ we have
\begin{eqnarray}
\alpha\circ\Phi\circ\mathcal{R}\circ\iota(b)=b\ .
\end{eqnarray}

There are two main frameworks of QECC.  In \textit{passive} error correction, our encoding map is such that any errors can only act trivially on our desired logical operators.  In \textit{active} error correction, error acts in a non-trivial way and we must find a recovery map $\mathcal{R}$ to ``undo" the errors for decoding.

\underline{Example} (0-form Global $\mathbb{Z}_2$ Symmetry):

Let $U$ implement the $\mathbb{Z}_2$ symmetry on the physical Hilbert space.
Define the error channel on operators by
\begin{equation}
\Phi_{\rm sym}(a)=U^\dagger a\,U\,,
\end{equation}
which matches Eq.~\eqref{Krausforms} with a single Kraus operator $V=U$.

\paragraph{Passive correction}
If the code is chosen so that the encoded algebra lands inside the invariant subalgebra,
\begin{equation}
\iota(\mB)\subset \mA^{\mathbb{Z}_2}:=\{a\in\mA~|~U^\dagger a\,U=a\}\,,
\end{equation}
then $\Phi_{\rm sym}$ acts trivially on all encoded observables:
$\Phi_{\rm sym}(\iota(b))=\iota(b)$ for every $b\in\mB$.
Hence the identity recovery suffices and
\begin{equation}
\alpha\circ \Phi_{\rm sym}\circ \iota=\alpha\circ \iota=\mathrm{id}_{\mB}\,,
\end{equation}
i.e. exact passive error correction.

\paragraph{Active correction}
Suppose we also consider non-observable\footnote{In AQFT, the observable algebra is the invariant subalgebra; charged fields connect sectors and are not observables.} charged fields that connect superselection sectors, e.g. an odd operator $\phi$ with $U^\dagger \phi\,U=-\phi$.
Since $\Phi_{\rm sym}$ is a unitary automorphism, these errors are still exactly correctable by choosing the inverse automorphism as recovery,
\begin{equation}
\mathcal{R}_{\rm sym}(a)=U\,a\,U^\dagger\,,
\end{equation}
so that for all $b\in\mB$,
\begin{equation}
\alpha\circ \Phi_{\rm sym}\circ \mathcal{R}_{\rm sym}\circ \iota(b)
=\alpha\big(\mathrm{id}_\mA(\iota(b))\big)=b\,.
\end{equation}

\underline{Example} (The Stinespring dilation of Kramers Wannier symmetry):

If we take the action of the symmetry operators to be the errors, in the Stinespring dilated Hilbert space described in Appendix~\ref{sec:LatticeKW}, the KW defects commuting past an operator insertion causes the error.  However, the defect “movement’’ gates $R_{D,\ell}$ act unitarily:
\begin{equation}
U_{\rm seq}:=\prod_{\ell} R_{D,\ell}\,,\qquad
\Phi_{\rm seq}(a)=U_{\rm seq}^\dagger a\,U_{\rm seq}\,.
\end{equation}
Thus, if we take the starting point as the Hilbert space with two domain walls already present, we are in the setting of invertible maps.  We again have a straightforward notion of passive and active exact QEC:
\begin{align}
\mathcal{C}_{\rm passive}&=\{\,b\in\mB~|~[\iota(b),U_{\rm seq}]=0\,\}\,,\\
\mathcal{R}_{\rm seq}(a)&=U_{\rm seq}\,a\,U_{\rm seq}^\dagger
\end{align}
In this case, $\mathcal{C}_{\rm passive}$ only contains the identity and $\eta$.

\subsubsection{Non-invertible Symmetry and the Petz Recovery Map}

Now let us return to the setting of the original, defect free Hilbert space of the Krammers Wannier example.  Here, we do not have a simple Kraus decomposition or unitary representation of the non-invertible symmetry action.  Instead, we will utilize the Petz recovery map as below.

Let $\rho$ be any faithful normal state on $\mA$ and define $\rho\circ\Phi_{\text{KW}}$ on $\mA$ by $(\rho\circ\Phi_{\text{KW}})(a):=\rho(\Phi_{\text{KW}}(a))$ as in \cite[Thm.\,1; App.\,B.4]{furuya2022real}. The Petz dual (recovery) of $\Phi_{\text{KW}}$ relative to $\rho$ is the CP map
\begin{equation}\label{eq:Petz-Heis}
\mathcal{R}^{\text{Petz}}_{\rho}(b)
\;=\;
\rho^{-1/2}\,\Phi_{\text{KW}}^{*}\!\big(\rho^{1/2}\,b\,\rho^{1/2}\big)\,\rho^{-1/2},
\qquad b\in\mA,
\end{equation}
where $\Phi_{\text{KW}}^{*}$ is the trace-dual (Schrödinger) channel. In a Kraus form for $\Phi_{\text{KW}}$,
\begin{equation}
\Phi_{\text{KW}}(a)=\sum_{r} V_r^\dagger a V_r
\quad\Longrightarrow\quad
\Phi_{\text{KW}}^{*}(\sigma)=\sum_{r} V_r\,\sigma\,V_r^\dagger,
\end{equation}
so that
\begin{equation}\label{eq:Petz-Kraus}
\mathcal{R}^{\text{Petz}}_{\rho}(b)
\;=\;
\sum_{r}\,
\rho^{-1/2}\,V_r\,\rho^{1/2}\,b\,\rho^{1/2}\,V_r^\dagger\,\rho^{-1/2}.
\end{equation}

Since we are restricting ourselves to the original Hilbert space and algebra, we cannot work with the form of \ref{eq:Petz-Kraus} which involves the enlarging isometries.  Instead, we work with the dual form directly.  Its action on the algebra can be deduced from the requirement that $\braket{a|\Phi_\text{KW}(b)} = \braket{\Phi^*_{\text{KW}}(a)|b}$, or equivalently $\text{tr}( a\Phi_\text{KW}(b)) = \text{tr}(\Phi^*_{\text{KW}}(a)b)$.  The resulting actions are

% \color{blue}

% \begin{eqnarray}
%     &&\Tr(\iota_1(X_{j})Z_j'Z_{j'+1}) = \Tr(X_{j}\iota^*_1(Z_{j'}Z_{j'+1})\\&&\Tr(\iota_1(X_{j})Z_j'Z_{j'+1}) = \Tr(X_{j}\iota^*_1(Z_{j'}Z_{j'+1})\\&&\delta_{j,j'} = \Tr(X_{j}\iota^*_1(Z_{j'}Z_{j'+1})\\
%     &&\implies \iota_2^*(Z_{j'}Z_{j'+1}) = 2X_{j'}
% \end{eqnarray}

% \begin{eqnarray}
%     &&\Tr(\iota_2(Z_jZ_{j+1})X_{j'}) = \Tr(Z_jZ_{j+1}\iota^*_2(X_{j'}))\\&&\Tr(X_{j+1}X_{j'}) = \Tr(Z_jZ_{j+1}\iota^*_2(X_{j'}))\\&&\delta_{j+1,j'} = \Tr(Z_jZ_{j+1}\iota^*_2(X_{j'}))\\
%     &&\implies \iota_2^*(X_{j'}) = \frac{1}{2}Z_{j'-1}Z_{j'}
% \end{eqnarray}

% \color{black}

\begin{eqnarray}
\Phi^*_\text{KW}(X_j) = \frac{1}{2}Z_{j-1}Z_{j}\\
\Phi^*_\text{KW}(Z_jZ_{j+1}) = 2X_j
\end{eqnarray}

Then Eq.~\eqref{eq:Petz-Heis} gives explicit formulas once a faithful reference state $\rho$ is fixed.
On a finite chain with periodic boundary conditions, we can take $\rho$ to be the maximally mixed state. Since $\rho^{1/2}$ is proportional to the identity, \eqref{eq:Petz-Heis} reduces to
\[
\mathcal{R}^{\text{Petz}}_{\tau}(b)=\Phi_{\text{KW}}^{*}(b).
\]

Hence the composition on the generators of $\mathcal{C}$ act as identity up to rescaling,
\begin{equation}\label{eq:tau-Petz-compose}
(\mathcal{R}^{\text{Petz}}_{\tau}\!\circ\Phi_{\text{KW}})(X_j)=2\,X_j,\qquad
(\mathcal{R}^{\text{Petz}}_{\tau}\!\circ\Phi_{\text{KW}})(Z_jZ_{j+1})
=\tfrac{1}{2}\,Z_jZ_{j+1}
\end{equation}

Note that the operator
$\textrm{KW}^\dagger$ has the desired recovery property.  It acts naturally on the dual space, and we have that $(\text{KW})(\text{KW}^\dagger) = (\text{KW})(\text{KW})T^{-1} = (1+\eta)TT^{-1} = (1+\eta)$, so on the $\mathbb{Z}_2$ even operators of interest the map acts \color{blue}proportional to \color{black} identity.  Thus, for the KW example, a representation of the Petz recovery map is the conjugate transpose of the non-invertible symmetry operator, which is the dual map in this canonical spin chain representation.

The lack of a passively correctable subalgebra in the case of lattice KW can be related to the Haag duality violation of the symmetric subalgebra \cite{shao2025additivityhaagdualitynoninvertible}.  In Ref.~\cite{shao2025additivityhaagdualitynoninvertible}, the only example of a duality violating subalgebra for topologically trivial regions was KW\footnote{This is in contrast to continuum examples for we know that Haag duality must hold for topologically trivially regions, unless the a symmetry is spontaneously broken\cite{halvorson2006algebraicquantumfieldtheory}}.  This is because of the mixing of spacetime and internal symmetries, which is also the reason that there are no local, invariant operators under the action of the symmetry map.  This mixing with translation means that the symmetry algebra does not categorify.  In later examples, we will see that symmetry algebras having a categorification automatically yields a conditional expectation onto some non-trivial passively correctable subalgebra \cite{wang2025manyfacesofnis}.

% \subsubsection{GNS Representation and QEC}

\begin{table}[t]
\caption{Summary of error correction maps and implementations for KW}
\label{tab:kw-summary}
\centering
\scriptsize
\setlength{\tabcolsep}{4pt}\renewcommand{\arraystretch}{1.05}
\begin{tabularx}{\textwidth}{lXX}
\toprule
\textbf{Representation} & \textbf{Noise map} & \textbf{Recovery map} \\
\midrule
Periodic lattice (N site) &
\textbf{Non-invertible} operator $\textrm{KW}$ &
\textbf{Non-invertible} dual operator $\textrm{KW}^\dagger$ \\
Ancilla lattice (Stinespring) &
\textbf{Unitary} defect movement
$\displaystyle \mathcal N(O)=U^\dagger OU$,
$\ U=\prod_j R_{d\ell}(j)& \textbf{Unitary}
conjugate map UOU^\dagger$ \\
GNS Hilbert space &
\textbf{Non-unitary} GNS implementer $D$  & \textbf{Unitary} GNS Petz recovery implementer \\
\bottomrule
\end{tabularx}
\end{table}

\subsection{Continuum Examples}

Building on the lattice example we saw earlier, we will frame well-known examples of non-invertible symmetries in continuum quantum field theory in terms of error correcting codes.  To this end, it is helpful to use the DHR approach to algebraic quantun field theory (AQFT), which we review here.

First, we say that a local quantum theory $\mathcal{T}$ on spacetime manifold $M$ assigns an operator algebra $\mA$ to regions of spacetime $R\in M$.   In the DHR approach, one considers ball-localized, gauge invariant, transportable \textit{observable} operators as the fundamental objects of interest.  We can reconstruct the gauge group, superselection sectors, and other data of interest using properties of the algebra of observables (see eg. Ref.~\cite{haag2012local, halvorson2006algebraicquantumfieldtheory} for a review).  \SA{Are field operators unobservable because charge is ball-localized or because they aren't gauge invariant?}  Field operators are \textit{unobservable} because their existence cannot be detected be an observer in a spacelike complimentary region (this excludes i.e. electric charges whose presence can be detected due to Gauss' Law).   We will denote the algebra of all operators as $\mA_\mathcal{F}$ and the algebra of the observable operators as $\mathcal{A}_\mathcal{O}$.  Non-local objects such as extended symmetry defects or Wilson lines that do not end are not natural concepts in this framework, so care must be taken in translating generalized symmetry examples to the DHR framework\footnote{The BF approach, which we will again mention in Section~\ref{sec:tqft}, attempts to circumvent these challenges by considering `cone-localized' charges instead of ball-localized}.

\underline{Example} (0-form $\mathbb{Z}_2$ symmetry reconstruction\footnote{We omit the word ``global" here because from the perspective of AQFT, we do not have access to global data about the theory.  However, as with the usual notion of global symmetry. we assume that the transformation acts uniformly on a fixed time slice.}):

Here we demonstrate how we can diagnose the existence of a 0-form $\mathbb{Z}_2$ gauge group from properties of local algebras \cite{CasiniEntropicOrderParameters_2021}.  Consider a theory 1+1d theory $\mathcal{T}$ on a topologically trivial region $\mathcal{R}$.  We further define two disjoint, ball-like subregions $\mathcal{R}_1$ and $\mathcal{R}_2$ in $\mathcal{R}$.  Thus, $\mathcal{R}/(\mathcal{R}_1\cup\mathcal{R}_2)$ is a region with non-trivial 0-th homotopy group $\pi_0$ as in the Figure~\ref{fig:z2intertwiner}.

\begin{figure}[t]
\centering
\begin{tikzpicture}[x=1cm,y=1cm, every node/.style={font=\small}]

% Big region R (a fixed-time slice picture)
\draw[line width=1.2pt] (0,0) -- (10,0);
\node[above] at (5,0.20) {$\mathcal{R}$};

% Complement components (gray)
\draw[line width=5pt, color=gray!35] (0,0) -- (2,0);
\draw[line width=5pt, color=gray!35] (3.6,0) -- (6.4,0);
\draw[line width=5pt, color=gray!35] (8,0) -- (10,0);

\node[below, color=gray!70!black] at (1, -0.28) {$C_1$};
\node[below, color=gray!70!black] at (5, -0.28) {$C_2$};
\node[below, color=gray!70!black] at (9, -0.28) {$C_3$};

% Subregions R1 and R2 (blue)
\draw[line width=5pt, color=blue!55] (2,0) -- (3.6,0);
\draw[line width=5pt, color=blue!55] (6.4,0) -- (8,0);

\node[above, color=blue!70!black] at (2.8,0.40) {$\mathcal{R}_1$};
\node[above, color=blue!70!black] at (7.2,0.40) {$\mathcal{R}_2$};

% Midpoints inside R1 and R2 for the intertwiner endpoints
\coordinate (midR1) at (2.8,0);
\coordinate (midR2) at (7.2,0);

% Mark endpoints (optional)
\fill (midR1) circle (2pt);
\fill (midR2) circle (2pt);

% Bi-local intertwiner connecting the midpoints
\draw[dashed, thick, bend left=35]
(midR1) to node[above] {$\mathcal{I}=\psi_1\psi_2$} (midR2);

% Text indicating disconnected complement (nontrivial pi_0)
\node[below] at (5,-0.90)
{$\mathcal{R}\setminus(\mathcal{R}_1\cup \mathcal{R}_2)=C_1\sqcup C_2\sqcup C_3
\quad\Rightarrow\quad \pi_0\neq 0$};

\end{tikzpicture}
\caption{A fixed-time slice in 1+1d with two disjoint ball-like regions $\mathcal{R}_1,\mathcal{R}_2\subset\mathcal{R}$. The complement inside $\mathcal{R}$ is disconnected (nontrivial $\pi_0$), and one can visualize bi-local intertwiners $\mathcal{I}=\psi_1\psi_2$ generated in $\mathcal{R}_1$ and $\mathcal{R}_2$.}
\label{fig:z2intertwiner}
\end{figure}

The observable operator content of $\mathcal{R}_1$
and $\mathcal{R}_2$ separately are $\mathbb{Z}_2$ even operators $\sigma$.  However, we can also have bi-locally generated operators that are charge neutral overall, but that are products of $\mathbb{Z}_2$ odd unobservables $\psi$.  We call these objects $\textit{intertwiners}$, $\mathcal{I} = \psi_1\psi_2$ .
Their existence corresponds to violations of \textit{additivity}, meaning $\mA(\mathcal{R}_1 \vee \mathcal{R}_2) \neq \mA(\mathcal{R}_1) \vee \mA(\mathcal{R}_2)$.  Hence we say that a violation of additivity signals the existence of a 0-form symmetry for regions with non-trivial $\pi_{d-2}$.  It is important to note that even for a (discrete) global gauge group $G$, restricting to the gauge invariant sector of the theory leads to the same violation of identity.  The gauge invariant sector of a theory is clearly a physical theory, so a violation of additivity does not automatically mean that the theory is invalid.

\subsubsection{0-form Rep($G$) symmetry}
\label{sec:repg}
In continuum quantum field theory, a straightforward construction of 0-form non-invertible symmetry comes from gauging an invertible 0-form symmetry with finite, non-Abelian group $G$ in 1+1d \cite{bhardwaj2018gaugingfinitesymmetries}.  The resulting gauged theory will contain extended co-deminsion 1 operators that are Wilson loops for the gauge field.  These operators are labeled by irreducible representations of $G$, so when the group is non-Abelian, the fusion of these operators is in general non-invertible.  We can equivalently consider the defect lines corresponding to Wilson loops as implementing the symmetry.  The collection of these defect lines along with physical consistency conditions form the data of a unitary fusion category $\mathcal{C} = \text{Rep}(G)$.

As noted in~\cite{harlow2025disjointadditivitylocalquantum}, gauging a finite group symmetry is different from simply projecting onto the subspace of symmetry invariant operators. The former should be viewed as a discrete Fourier transform that exchanges information about the charged and twisted sectors of the gauged and ungauged theories, while the latter is a process that does not result in a complete theory with a dual symmetry~\cite{bhardwaj2018gaugingfinitesymmetries, bhardwaj2023lecturesgeneralizedsymmetries}.

Here we will review the above facts in operator algebraic language.  We start with an algebra of observables $\mA$. Assume that $\mA$ has some internal symmetry specified by automorphisms of $\mA$
\begin{eqnarray}
\{\alpha_g\},\;g\in G,
\end{eqnarray} where $G$ is a discrete non-Abelian group.  We further require that the algebra $\mA$ satisfies \textit{completeness}, meaning Haag duality holds for arbitrary regions.  As such there is no group $H$ such that $\mA = \mathcal{F}/H$, where $\mathcal{F}$ is some parent theory containing unobservable, charge creating fields.

\begin{figure}[t]
\centering

% ---------------- Fig 2a ----------------
\begin{subfigure}{\linewidth}
\centering
\begin{tikzpicture}[x=1cm,y=1cm, every node/.style={font=\small}]

% Big region R (fixed-time slice)
\draw[line width=1.2pt] (0,0) -- (10,0);
\node[above] at (5,0.20) {$\mathcal{R}$};

% Complement pieces (gray)
\draw[line width=5pt, color=gray!35] (0,0) -- (2,0);
\draw[line width=5pt, color=gray!35] (3.6,0) -- (6.4,0);
\draw[line width=5pt, color=gray!35] (8,0) -- (10,0);

% Subregions R1 and R2 (blue)
\draw[line width=5pt, color=blue!55] (2,0) -- (3.6,0);
\draw[line width=5pt, color=blue!55] (6.4,0) -- (8,0);

\node[above, color=blue!70!black] at (2.8,0.42) {$\mathcal{R}_1$};
\node[above, color=blue!70!black] at (7.2,0.42) {$\mathcal{R}_2$};

% Mark the boundaries of R1 and R2
\draw[densely dashed] (2,0.32) -- (2,-0.32);
\draw[densely dashed] (3.6,0.32) -- (3.6,-0.32);
\draw[densely dashed] (6.4,0.32) -- (6.4,-0.32);
\draw[densely dashed] (8,0.32) -- (8,-0.32);

% Twist operator: centered in R1, but "smeared" slightly outside its boundaries
\coordinate (midR1) at (2.8,0);

% soft smear (stacked translucent blobs so it leaks into the complement)
\fill[orange!80, opacity=0.10] (midR1) ellipse (1.35 and 0.40); % widest: extends beyond R1
\fill[orange!80, opacity=0.16] (midR1) ellipse (1.05 and 0.34);
\fill[orange!80, opacity=0.24] (midR1) ellipse (0.78 and 0.28);

% core marker + label (label color matches smearing color)
\fill[orange!85!black] (midR1) circle (1.7pt);
\node[above, text=orange!80] at (2.8,0.80) {$\tau_c$};

% Identity text above the R2 label, aligned in the same "band" as tau_c above R1
\node[above, text=orange!80] at (7.2,0.80) {\scriptsize identity on $\mathcal{R}_2$};

\end{tikzpicture}
\subcaption{Twist $\tau_c$ for two disjoint regions in 1+1d. $\tau_c$ implements a charge measurement in $\mathcal{R}_1$ while acting trivially on $\mathcal{R}_2$. The smeared support crosses the boundary of $\mathcal{R}_1$ so that $\tau_c$ is not strictly localized to $\mathcal{R}_1$ alone.}
\label{fig:2dtwist:a}
\end{subfigure}

\vspace{0.8em}

% ---------------- Fig 2b ----------------
\begin{subfigure}{\linewidth}
\centering
\begin{tikzpicture}[x=1cm,y=1cm, every node/.style={font=\small}]

% Big region R (fixed-time slice)
\draw[line width=1.2pt] (0,0) -- (10,0);
\node[above] at (5,0.20) {$\mathcal{R}$};

% Complement pieces (gray)
\draw[line width=5pt, color=gray!35] (0,0) -- (2,0);
\draw[line width=5pt, color=gray!35] (3.6,0) -- (6.4,0);
\draw[line width=5pt, color=gray!35] (8,0) -- (10,0);

% Subregions R1 and R2 (blue)
\draw[line width=5pt, color=blue!55] (2,0) -- (3.6,0);
\draw[line width=5pt, color=blue!55] (6.4,0) -- (8,0);

\node[above, color=blue!70!black] at (2.8,0.42) {$\mathcal{R}_1$};
\node[above, color=blue!70!black] at (7.2,0.42) {$\mathcal{R}_2$};

% Mark the boundaries of R1 and R2
\draw[densely dashed] (2,0.32) -- (2,-0.32);
\draw[densely dashed] (3.6,0.32) -- (3.6,-0.32);
\draw[densely dashed] (6.4,0.32) -- (6.4,-0.32);
\draw[densely dashed] (8,0.32) -- (8,-0.32);

% Boundary-localized depiction: epsilon-neighborhoods around endpoints of R1
\coordinate (bL) at (2,0);
\coordinate (bR) at (3.6,0);

% soft smear around left boundary point (straddles R1 and complement)
\fill[orange!80, opacity=0.10] (bL) ellipse (0.45 and 0.40);
\fill[orange!80, opacity=0.16] (bL) ellipse (0.35 and 0.34);
\fill[orange!80, opacity=0.24] (bL) ellipse (0.25 and 0.28);

% soft smear around right boundary point
\fill[orange!80, opacity=0.10] (bR) ellipse (0.45 and 0.40);
\fill[orange!80, opacity=0.16] (bR) ellipse (0.35 and 0.34);
\fill[orange!80, opacity=0.24] (bR) ellipse (0.25 and 0.28);

% core markers at the boundary points
\fill[orange!85!black] (bL) circle (1.7pt);
\fill[orange!85!black] (bR) circle (1.7pt);

% optional epsilon brackets (visual cue)
\draw[<->, line width=0.7pt, color=orange!80] (1.78,-0.52) -- (2.22,-0.52)
  node[midway, below, text=orange!80] {\scriptsize $\epsilon$};
\draw[<->, line width=0.7pt, color=orange!80] (3.38,-0.52) -- (3.82,-0.52)
  node[midway, below, text=orange!80] {\scriptsize $\epsilon$};

% label indicating "boundary support"
\node[above, text=orange!80] at (2.8,0.80) {$\tau_c \ \text{near}\ \partial\mathcal{R}_1$};

% Identity text above the R2 label, aligned in the same band
\node[above, text=orange!80] at (7.2,0.80) {\scriptsize identity on $\mathcal{R}_2$};

\end{tikzpicture}
\subcaption{Boundary-localized picture of the twist.  The nontrivial content is encoded in an $\epsilon$-neighborhood of $\partial\mathcal{R}_1$, while $\tau_c$ commutes with operators strictly inside $\mathcal{R}_1$.}
\label{fig:2dtwist:b}
\end{subfigure}

\caption{Twist for two disjoint regions in 1+1d. Figure (a) depicts the twist and its defined support. Figure (b) depicts the equivalent viewpoint where the nontrivial action is localized near the boundary $\partial\mathcal{R}_1$.}
\label{fig:2dtwist}
\end{figure}



In the physics literature, gauging is defined such that if the original theory is complete, then the gauged theory is also complete.  In the setting of orbifolding 2d CFTs, this is the familiar notion of requiring the orbifolded theory to respect modular invariance.  This means that in addition to quotienting by the group action, we must include \textit{twisted sectors}, or \textit{twisted operators} which are invariant under the original symmetry action $\alpha_g$ but are charged under a dual symmetry of the gauged theory $\mA_{\text{gauged}}$.  In the operator algebraic approach, the objects that restore Haag duality for arbitrary regions are \textit{twists}~\cite{CasiniMagan2025ABJasU1symmetry}.  For two disjoint regions as considered in the intertwiner example, twists can be understood as implementing a charge measurement in one of the regions while doing nothing to the other region (Figure~\ref{fig:2dtwist}a).  As is the case with a generic charged operator in AQFT, each twist will correspond to a unitarily inequivalent representation, or superselection sector, of the operator algebra.  These are exactly the familiar twisted sectors from orbifolding.

By construction, twists are neutral under $\alpha_g$.  In the non-Abelian case, we have to take sums within a specific conjugacy class $c$ so that the twists are charge neutral under the symmetry action overall~\cite{CasiniEntropicOrderParameters_2021}:

\begin{eqnarray}
U_h\tau_gU_h^\dagger = \tau_{hgh}\\
\tau_c = \sum_{g\in c}\tau_g
\end{eqnarray}

We see that for non-Abelian $G$, charge neutral twists fuse non-invertibly like the group's irreducible representations, forming a Rep($G$) fusion rule.  Thus, we see that that the gauged theory $\mA_\text{gauged} = \mA/G\;\vee\;\{\tau_c\}$ has a non-invertible Rep($G$) symmetry whose charged objects are the twists $\tau_c$.  We call this algebra a \textit{field algebra} as the twists are fields for the dual symmetry of the gauged theory.

One may notice an apparent mismatch that the charged operators under a $0$-form symmetry should be $0$-dimensional objects, whereas the twists are supported on finite $1$-dimensional regions.  This can be reconciled by understanding that the non-trivial information of the twist is encoded at the boundary of the supported region $\partial R_1$.  In the gauged theory, the twists commute with all operators inside $R_1$.  The twists also live in the commutant of charge neutral observables in $R_1' \cup R_2'$.  Thus, we claim that the non-trivial information of the twists in the gauged theory is encoded in the action of the symmetry on the $\epsilon$ width smearing around the boundary, which can effectively be taken to be $\partial R_1$ (Figure~\ref{fig:2dtwist}b).  As we take the limit $\epsilon \to 0$, the charged objects can be understood as 0-form operators localized at the endpoints of the region as expected.

\SA{Here I think we can remark about previous works and the crossed product structure that arises from gauging.  I think it would be good to cite the work from Marc \cite{klinger2025theory} and Muger \cite{Muger2005conformalorbifold, muger1999soliton}.  Marc discusses the symmetry broken case and Muger covers orbifolding CFTs for finite groups.  See especially section 2.2 of Muger's work \cite{muger1999soliton}.  He extends the invariant subalgebra by disorder operators that he defines which implement the symmetry action in a half-region, which is very similar to the "include twists after projecting" picture we discuss.  The structure he arrives at in Eq. 2.10 and Lemma 2.2 are the crossed product structure we've been discussing, and he remarks at the bottom of that page about the dual symmetry that arises via Takesaki duality.  I think the difference in our case is that our groups are non-Abelian, which is why we take the crossed product with $\hat{G}$ instead of $G$}

As clear from the above discussion, the algebra of observables in the gauged theory will be decomposed into twisted superselection sectors labeled by the conjugacy classes of the gauge group $G$.  We also have a set of superoperators for the symmetry:

\begin{eqnarray}
&\mA = \bigoplus_{[g] \in G} \mA_{[g]}\\
&\{\Phi_{W_c}| \text{$c$ an irrep. of $G$}\}
\end{eqnarray}

The action of the superoperators simply picks up a factor given by the group character depending on the twisted sector of the operator in the ungauged theory $$\Phi_{W_c}(a_{[g]}) = \chi_c([g])a_{[g]}$$

In the context of the Lagrangian formulation of QFT, this corresponds to Wilson operators acting on an a closed loop with non-trivial holonomy. We would like to first construct a conditional expectation to map the full algebra to the passively correctable subalgebra $\mA_{\text{passive}}$.  Define the map $E_{\mC}:\mA_{\text{gauged}}\to\mA_{\text{passive}}$ by

\begin{equation}\label{eq:E-def-RepG}
E_{\mC}(a_g)=\frac{1}{D^2}\sum_{c\in C}d_c\Phi_{W_c}(a_{[g]}).
\end{equation}

Here, $D$ is the total quantum dimension of $\mC$ and $d_c$ is the quantum dimension of the simple object $c$ in the category (see Appendix \ref{sec:condexp} for a review of category theory terminology in this context).  When the category is Rep($G$), the quantum dimensions of the object are the dimension of the irrep.  We prove in Appendix~\ref{sec:condexprepg} that this map satisfies the definition of a conditional expectation and acts in the desired way on the trivial and actively-correctible subalgebras.

As in the KW example, active error correction corresponds to acting with the symmetry defect dual to whichever non-invertible symmetry error occured.  Concretely, if an error corresponds to a simple defect $c\in\mathcal{C}$ acting on a region, one applies $c^*$ and measures the fusion outcome in
\begin{equation}\label{eq:fusion-measurement-general}
c\otimes c^* \;=\;\bigoplus_{a\in\mathcal{C}} N_{cc^*}^{\;\;a}\,a.
\end{equation}
In a topological theory (or in the topological sector of a CFT), the probability to obtain outcome $y$
is proportional to its quantum dimension:
\begin{equation}\label{eq:fusion-prob}
p(a\,|\,c\otimes c^*)=\frac{N_{cc^*}^{\;\;a}\,d_a}{d_c\,d_{c^*}}=\frac{N_{cc^*}^{\;\;a}\,d_a}{d_c^{\,2}}.
\end{equation}
If the measured outcome $a$ is invertible, the remaining correction is deterministic via applying $a^*$.
If the measured outcome is again non-invertible, one repeats.

\underline{Example} (0-form Rep($S_3$) Symmetry):

The simplest example of the above construction comes from gauging a theory with a $G = S_3$ symmetry in 1+1$d$~\cite{bhardwaj2025gappedphaseswithnoninvertiblesymm}.  We denote the ungauged theory by $\mathcal{T}'$ and the gauged theory by $\mathcal{T} = \mathcal{T'}/G$.  There are 3 irreducible representations for $S_3$, which we denote by  $1$, $\epsilon$, and $\sigma$, that fuse as $\epsilon \times \epsilon = 1$, $\epsilon \times \sigma = \sigma$, and $\sigma \times \sigma = 1 + \sigma + \epsilon$.  In $\mathcal{T}'$, the twisted sectors are labeled by the conjugacy classes, which are transpositions $\tau$ or 3-cycles $c$.  The character table is given in Table \ref{tab:s3chartable}.

\begin{table}[h]
\centering
\begin{tabular}{lccc}
\hline
irrep $R$ & $\chi_R(e)$ & $\chi_R(\tau)$ & $\chi_R(c)$ \\
\hline
$1$ & $1$ & $1$ & $1$ \\
$\varepsilon$ & $1$ & $-1$ & $1$ \\
$\sigma$ (2d) & $2$ & $0$ & $-1$ \\
\hline
\end{tabular}
\caption{Character table of $S_3$}
\label{tab:s3chartable}
\end{table}

Eq. \eqref{eq:E-def-RepG} then gives:

\begin{eqnarray}
&E_C(a_e) &= \frac{1}{6}(1(1) + 1(1) +2(2))a_e \\
&&=a_e\\
&E_C(a_\tau) &= \frac{1}{6}(1(1) - 1(1) +2(0))a_\tau \\
&&=0\\
&E_C(a_c) &= \frac{1}{6}(1(1) + 1(1) +2(-1))a_c \\
&&=0
\end{eqnarray}

as expected.

For active error correction, the only non-invertible symmetry is implemented by $\Phi_\sigma$ which is a self dual object.  Then, if we act with $\Phi_\sigma$ again on the local operator, with probability $\frac{1}{2}$ we will get a fusion outcome that is deterministically correctable (1 or $\epsilon$).  Alternatively, with probability $\frac{1}{3}$, we will get $\sigma$ again and have to repeat the fusion.  Generalizing, we see that with $n$ fusion outcome measurements we will have deterministic error correction with probability $1-(\frac{1}{3})^n$.

\subsubsection{Higher-form Rep($G$) Symmetry}

We can generalize the procedure of the previous section to generate higher form non-invertible symmetries of the Rep($G$) type via gauging.  Starting with an algebra of observables $\mA$ in $d$ spacetime dimensions with a $0$-form non-Abelian symmetry $G$, we quotient out by the operators that are not invariant under the automorphism action and add twists labeled by conjugacy classes $c = [g]$.  The charged objects under the dual symmetry can be thought of as the $d-2$ dimensional boundaries of the $d-1$ dimensional regions that the twists are supported in.  We recover the familiar understanding that gauging a $0$-form non-Abelian $G$ symmetry in $d$ spacetime dimensions gives a theory with $(d-2)$-form Rep($G$) non-invertible symmetry.

The error correction code that corresponds to the gauged theory has the same structure in general dimensions.  The conditional expectation in Eq.~\ref{eq:E-def-RepG} mapping to the passively correctable subalgebra and the probabilities for active error correction in Eq.~\ref{eq:fusion-prob} depend only on the abstract categorical data, so the previous analysis carries over.



\subsubsection{Verlinde Lines in 2d CFT}

Another canonical example of non-invertible 0-form symmetries in quantum field theory is topological defect lines in 2d CFT~\cite{shao2019tdlrgflow, shao2024whatsundonetasilectures}.  For CFTs, projectors onto the trivial representation of operators has been explored in the literature due to interest in symmetry resolved entanglement entropy \cite{tachikawa2023asymptoticdos,bastidadas2024catsree}.  Here we review those results and rephrase them in our language.  We will assume periodic boundary conditions for simplicity, noting that care must be taken as the conditional expectation would need to be adjusted depending on the choice of boundary conditions \cite{heymann2025revisitingsrefornis}.

Many familiar concepts in $2d$ CFTs can be recovered from the DHR approach~\cite{Kawahigashiclassificationlocalconformalnets}. Starting with the algebra of the stress tensor, or the identity block, along with conformal symmetry conditions, one can reconstruct the conformal families and the resulting superselection sectors from acting with the primaries in the vacuum representation.  For rational conformal field theories, the local operators can be organized in terms of holomorphic and antiholomorphic representations of the chiral algebra $a\in\mA_\text{chiral}$.

\begin{eqnarray}
&\mathcal{A} = \bigoplus_a\mathcal{A}_a\otimes \bar{\mathcal{A}}_{a^*}
\end{eqnarray}

An important property of CFTs is modular invariance, where the torus partition function is invariant under the modular group $SL(2, \mathbb{Z})$.  Ref.~\cite{casinietal2024modular} argues that modular invariance for CFTs is equivalent to completeness of the algebras from an AQFT perspective.  Many derivations in $2d$ CFT rely on modular invariance as a constraint, so physicists often work with the enlarged field algebra rather than just the algebra of charge neutral local operators. \SA{Maybe some editorial comment here about lack of completeness, or additivity or duality violations not implying that anything is wrong with the theory since it depends on choices or context?}  

Starting with the complete algebra of a diagonal rational CFT, one can derive the existence of generalized twists known as \textit{Verlinde lines} that are in a one-to-one correspondence with local primary operators \cite{petkova2001generalised}.  These twists commute with the generators of the Virosaro algebra by construction, so they are topological in the path integral sense.  Algebraically, they can be understood as intertwining representations with different boundary conditions.  Importantly. their action on the Hilbert space need not be unitarily implemented, in which case they are examples of non-invertible symmetries.  Their action on local operators is fixed by modular invariance to be

\begin{eqnarray}
\Phi_{\mathcal{L}_a}(O_b) = \frac{S_{ab}}{S_{1b}} O_b,
\end{eqnarray}

where $a$ is the label of the Verlinde line (or equivalently the label of a primary field), $b$ is the representation of the local operator and $S_{ab}$ is an element of the modular S-matrix.

The full set of superoperators correspond to the Verlinde lines of the theory along with the symmetry automorphisms.  The these objects and their associated fusion rules form the mathematical structure of a fusion category $\mC$.

\begin{eqnarray}
&\{\Phi_{\mathcal{L}_{c}}| c \in \mC\}
\end{eqnarray}

In some cases, a subset of the defect lines generates a particular fusion category symmetry $\mathcal{F}$, as we will see in an example below. 

We can use the Hilbert space projectors in Ref.~\cite{tachikawa2023asymptoticdos} to construct our conditional expectation since for CFTs there is an exact state-operator correspondence:

\begin{equation}\label{eq:E-def-cft}
E_{C}(O_b)=\frac{1}{D_\mathcal{C}^2}\sum_{c\in \mathcal{F}}d_c\Phi_{\mathcal{L}_c}(O_b)
\end{equation}

The quantum dimension in terms of the S-matrix is $d_c = \frac{S_{1c}}{S_{11}}$, with $D_\mathcal{C} = \sum_{c\in\mathcal{F}} d_c^2$.  When $\mathcal{F}$ is the full set of Verlinde lines, the conditional expectation returns the coefficient of the identity primary.  We show in Appendix \ref{sec:condexpcft} that the above is a conditional expectation onto the $\mathcal{F}$-invariant subalgebra.  As in the Rep($G$) case, active error correction is implemented by the superoperator dual to whichever error occurred, with the quantum dimensions and fusion coefficients determining the probability of successful correction.

\underline{Example} (KW Symmetry in the Ising CFT):

This is the simplest example of non-invertible symmetry in quantum field theory and can be understood by taking the continuum limit of the spin chain in the beginning of Section \ref{sec:IntrotoNUCNIS}.  The operators $Z_l$ and $Z_lZ_{l+1} - X_l$ flow to primaries $\epsilon$ and $\sigma$ respectively.  The primaries satisfy fusion rules $\epsilon \times \epsilon = 1,\;\epsilon \times \sigma = \epsilon,\;\sigma \times \sigma = 1 + \epsilon$. In the continuum, lattice translation acts trivially on the low lying states, so we have $\mathcal{T} \rightarrow \mathbbm{1}$. The symmetry superoperator algebra becomes
\begin{eqnarray}
&\Phi_{\mathbb{Z}_2} \circ \Phi_{\mathbb{Z}_2} = \mathbbm{1}\\
&\Phi_{\mathbb{Z}_2} \circ \Phi_{\text{KW}} = \Phi_{\text{KW}} \circ \Phi_{\mathbb{Z}_2}\\
&\Phi_{\text{KW}} \circ \Phi_{\text{KW}} = \mathbbm{1} + \Phi_{\mathbb{Z}_2}
\end{eqnarray}

The modular S-matrix is
\begin{eqnarray}
S =
\begin{pmatrix}&1&1&\sqrt{2}\\&1&1&\sqrt{2}\\&\sqrt{2}&-\sqrt{2}&0
\end{pmatrix}
\end{eqnarray} in the $1, \epsilon, \sigma$ basis.
The conditional expectation \eqref{eq:E-def-cft} for the full Ising fusion-category symmetry gives
\begin{eqnarray}
E_{C}(O_1) &=& \frac{1}{4}\Big(1\cdot 1 + 1\cdot 1 + \sqrt2\cdot \sqrt2\Big)O_1 = O_1,\\
E_{C}(O_\epsilon) &=& \frac{1}{4}\Big(1\cdot 1 + 1\cdot 1 + \sqrt2\cdot (-\sqrt2)\Big)O_\epsilon = 0,\\
E_{C}(O_\sigma) &=& \frac{1}{4}\Big(1\cdot 1 + 1\cdot (-1) + \sqrt2\cdot 0\Big)O_\sigma = 0,
\end{eqnarray}
so the passively-correctable local subalgebra is the vacuum family as expected.

The only non-invertible defect is the self dual $\sigma$.  If an error corresponds to the action of $\sigma$
on a region, we apply $\sigma$ and measure the fusion outcome in $\sigma\times\sigma=1+\epsilon$.
In either case the remaining correction is deterministic since the fusion outcomes are each invertible.  Thus a single fusion-outcome measurement suffices for deterministic active correction.

\underline{Example} (Fibonacci Symmetry in the three-state Potts CFT):  In the three-state Potts CFT one can isolate a fusion subcategory $\mathcal{F} \simeq \text{Fib}$.  The projectors onto the symmetric subalgebra were considered in Ref.~\cite{heymann2025revisitingsrefornis}, with special attention given to boundary conditions for the purpose of computing entanglement entropy.

The theory has six primary fields, $1, \psi, \psi^
\dagger, \epsilon, \sigma, \sigma^\dagger$, but we can consider the subset of Verlinde lines $\{\mathcal{L}_1, \mathcal{L}_\epsilon\}$ to generate $\text{Fib}$.  They fuse as $\mathcal{L}_\epsilon \times \mathcal{L}_\epsilon = 1 + \epsilon$ with a corresponding quantum dimension $d_\epsilon = \varphi = \frac{1 + \sqrt{5}}{2}$.  The Verlinde line $\mathcal{L}_\epsilon$ acts in the vacuum representation on the conformal towers of $1, \psi$, and $\psi^\dagger$, and the conditional expectation reads:

\begin{eqnarray}
E_C = \frac{1}{1+\varphi^2}(\Phi_{\mathcal{L}_1}+\varphi \Phi_{\mathcal{L}_\epsilon})
\end{eqnarray}

One can verify that this acts in the desired way on the primary fields, projecting out the $\epsilon$, $\sigma$, and $\sigma^\dagger$ primaries and leaving the rest invariant.

If an error corresponds to $\tau$, we apply $\tau$ and measure the fusion outcome of
$\tau\times\tau=1+\tau$.  Using \eqref{eq:fusion-prob} with $d_\tau=\varphi$ gives
\begin{equation}
p(1)=\frac{1}{\varphi^2},\qquad p(\tau)=\frac{\varphi}{\varphi^2}=\frac{1}{\varphi}.
\end{equation}
After $n$ fusion-outcome measurements, the probability of having obtained $1$ at least once is
\begin{equation}
p(n)=1-\left(\frac{1}{\varphi}\right)^{n},
\end{equation}

\subsubsection{Anyons in 3d TQFT}
\label{sec:tqft}
\SA{Note: on the lattice it has been recently proven that quantum dimensions are integers - even more questions about constructing rigorous models for non-Abelian anyons with Rep($H$) statistics}
Finally, let us consider the non-invertible symmetry generated by anyon worldlines in topological quantum field theory (TQFT).  This example allows us to address tension in the literature between algebraic quantum field theory and non-invertible symmetries\color{brown}cite Magan, Shao Sorce, Harlow\color{black}.  The authors of \color{brown}cite Magan\color{black} have argued against the existence of intrinsic non-invertible symmetries on the basis of an axiomatic approach to algebraic quantum field theory for spacetime $d>2$.  Here, intrinsic means that the non-invertible symmetry does not originate from an invertible one as in Section \ref{sec:repg}.

On the other hand, anyonic particles exist in 2+1 spacetime dimensions with fractionalized charge and exotic braiding statistics.  The Wilson lines carried by the anyon are topological operators with non-invertible fusion rules which, according to modern terminology, qualifies them as non-invertible 1-form  symmetries \color{brown}cite McGreevy, Komorgodski\color{black}.  The resolution relies in assumptions of local structures in algebraic quantum field theory. Foremost, in Ref.~\cite{casini2025generalizationdhrtheoremhigher} the authors assumptions exclude strictly topological theories, which rules out the effective field theories describing anyons. Indeed, topological quantum field theories which have no local operators besides the identity. Second, they restrict to local algebras that are generated in ball-like regions.  They remark that
BF selection criteria, which localize charges in cones extending to  infinity, correspond to singular conical regions and can be analyzed with similar tools, but are not treated in the paper.  This distinction is physically important in $2+1$d because the standard DHR criterion automatically leads to symmetric statistics in $D>2$, whereas anyonic braid statistics arises for cone-localized charges~\cite{fredenhagen1989braid,frohlich1990braid, naaijkens2011stringlike,plaschkeAnyonsThesis}.

In this work we do not attempt to rigorously define topological field theories of anyons in the language of AQFT.  However, we can still use the language of completely positive maps and algebras of extended operators to discuss the passive and active error correction induced by anyon worldlines\footnote{This error correction picture for anyonic theories bears some similarity to paradigms in topological quantum computation.  Information stored in the fusion state space of the anyons, which is different from the setup we consider, but unwanted fusion and braiding of anyons corrupts this information.}.  For a $2+1$d topological order described by a unitary modular tensor category (MTC) $\mathcal{C}$,
the structure is exactly the same as in the Rep($G$) and RCFT examples : the relevant
superoperators are precisely the anyon types labeled by simple objects $c\in\mathcal{C}$, their eigenvalues on charge sectors
are fixed by modular data, and the conditional expectation is the same weighted
average as in Eqs.~\eqref{eq:E-def-RepG} and \eqref{eq:E-def-cft}.  The only essential difference is
that a TQFT has no nontrivial pointlike local algebra on a ball, so one works with an
extended operator algebra that detects 1-form data.

Concretely, let $\mA$ denote the annulus algebra which decomposes into charged sectors labeled by simple objects $a\in\mathcal{C}$,
\begin{equation}
\mA=\bigoplus_{a\in\mathcal{C}}\mA_a.
\end{equation}
An anyon loop of type $c$ wrapped along the nontrivial cycle defines a superoperator
$\Phi_{W_c}$ acting diagonally on $\mA_a$ with eigenvalue given by the modular S-matrix:
\begin{equation}\label{eq:anyons-loop-eigs}
\Phi_{W_c}(O_a)=\frac{S_{ca}}{S_{1a}}\,O_a,\qquad O_a\in\mA_a.
\end{equation}

Define the conditional expectation by the same weighted average as before:
\begin{equation}\label{eq:anyons-condexp}
E_{\mathcal{C}}(O_a)=\frac{1}{D^2}\sum_{c\in\mathcal{C}} d_c\,\Phi_{W_c}(O_a),\qquad
d_c=\frac{S_{1c}}{S_{11}},\qquad D^2=\sum_{c\in\mathcal{C}} d_c^2.
\end{equation}
Using \eqref{eq:anyons-loop-eigs} and the standard $S$-matrix orthogonality,
$E_{\mathcal{C}}$ projects onto the trivial-charge sector:
\begin{equation}
E_{\mathcal{C}}(O_a)=\delta_{a,1}\,O_a.
\end{equation}
The passively-correctable subalgebra is the vacuum sector of the extended
annulus algebra, as in the RCFT example.

Active correction is implemented in exactly the same way as
Eqs.~\eqref{eq:fusion-measurement-general}--\eqref{eq:fusion-prob}.  If an error changes the
topological charge sector to $x\in\mathcal{C}$, one applies $x^*$ and measures the fusion channel
in $x\otimes x^*=\bigoplus_a N_{xx^*}^{\;\;a}a$, with probabilities
$p(a|x\otimes x^*)=N_{xx^*}^{\;\;a}d_a/d_x^{\,2}$.  Whenever the outcome is invertible
($d_a=1$), the remaining correction is deterministic by applying $a^*$; otherwise one repeats.

\section{Discussion}

\color{brown}Future directions: Is our picture somehow dual to one where errors are protected by NIS? Can we think of i.e. a depolarizing channel composed with a unitary as a non-invertible symmetry, [our framework is ideal for studying microscopic descriptions of NIS in mixed state systems in both lattice and continuum settings]?  

On the lattice, there might be NIS that does not have a corresponding ring or categorification $\to$ could the error correction maps be something other than just the dual in this case?  For the examples considered in the continuum, we always have some categorification with a structure of duals, so the active error correction maps are always just the dual of the non-invertible symmetry superoperator in the categorical sense.

Deeper questions about locality and NIS, connection with QCAs

 Does the crossed-product understanding of gauging easily extend to non-invertible symmetries i.e. taking the product with respect to some simple object of a category rather than a group element action.
 
 Should mention that there are many other examples of NIS that involve more exotic mathematical descriptions, so our statements have not been proven in a general sense
 
 Note that we can always find a conditional expectation by passing to the WHA, although uniqueness is an issue \color{black}




\[\]
\pagebreak
\appendix

\section{Lattice Krammers Wannier Symmetry}
\label{sec:LatticeKW}


Consider a 1D spin chain on a lattice of length $L$. The operators $X_j$ for $j\in S=\{1,\cdots, L-1\}$ generate the Abelian group $\mathbb{Z}_2^{L-1}$. The group algebra associated with these operators is the Abelian algebra $\mA_S=\sum_{i=1}^{L-1} \alpha_j X_j$ with $\alpha_{j=1}\in\mathbb{C}$. We define the unitary
\begin{eqnarray}
U_{\mA^S}=\prod_{j=1}^{L-1}\lb \frac{1+i X_j}{\sqrt{2}}\rb\in \mA_S\ .
\end{eqnarray}
We would like to enlarge $\mA_S$ by adding  $X_L$, which we are going to treat differently because it will be a ``boundary" point. For our discussion, it is more convenient to think we are adding the nonlocal operator $\eta=X_1\cdots X_L$ (Of course, this is equivalent to adding $X_L$). We call the enlarged algebra $\mA$.

We repeat the same construction above, starting with $Z_jZ_{j+1}$ for $j\in S$ to define the Abelian algebra $\mB^S$ \color{brown} should add some note explaining that we choose to not add $Z_1Z_L$ to avoid multi-valued issue \color{black}.

% \color{green} and the unitary
% \begin{eqnarray}
% U_{\mB^S}=\prod_{j=1}^{L-1}    \lb \frac{1+i Z_jZ_{j+1}}{\sqrt{2}}\rb\in \mB_S
% \end{eqnarray}
% and enlarge it by including \sout{$\eta$}  $Z_{L-1}Z_{L} $ to obtain the Abelian algebra $\mB$.
% \color{black}

% \sout{The two algebras $\mA$ and $\mB$ are isomorphic, as can be established by the isomorphism}

We can establish a morphism from $\mA$ to $\mB$ as below:
\begin{eqnarray}
&&\iota_1:\mA\to \mathcal{B}:\qquad \iota_1(1)=1,\qquad \iota_1(\eta)= 1,\\
&&\forall j\in S:\quad\iota_1(X_j)=Z_jZ_{j+1}\\
&&\iota_1(X_1...X_{L-1})=Z_1Z_{L}\\
&&\iota_1(X_L)=\iota_1(\eta)\iota_1(X_1\cdots X_{L-1})=Z_LZ_1\ .
\end{eqnarray}

We also have a morphism from $\mB$ to $\mA$:

\begin{eqnarray}
&&\iota_2:\mB\to \mA:\qquad \iota_2(1)=1\\
&&\forall j\in S:\quad\iota_2(Z_jZ_{j+1})= X_{j+1}\\
&&\iota_2(Z_1Z_{L})= X_{2}...X_{L}
\end{eqnarray}
x
We note several interesting features.  First, the algebra $\mA$ is strictly larger than the algebra $\mB$ since it contains the operator $\eta$.  The full string of $Z_jZ_{j+1}$ in $\mB$ just 1.  As a consequence, we see that we can generate the projector $P_- = \frac{1-\eta}{2} \in \mA$.  \color{brown} $P_-$ in kernel \color{black}Also, note that $\iota_2$ sends $Z_jZ_{j+1}$ to $X_{j+1}$ instead of $X_j$ in order to preserve the commutation relation $[X_j,Z_{j+1}Z_{j+2}] = 0$ under the pair of morphisms.

% For $j$ and $j+s$ in $S$, the
% morphism $\iota_1$ sends the string $X_j\cdots X_{j+s}$ to $Z_jZ_{j+s+1}$ at the endpoints. As a result, we have
% \begin{eqnarray}
%     &&\iota_1\color{black}(X_1\cdots X_{L-1})=Z_1Z_L\\
%     &&\iota_1\color{black}(X_L)=\iota_1\color{black}(X_1\cdots X_{L-1})\iota(\eta)=Z_1Z_L
% \end{eqnarray}

Since $\iota_1$ and $\iota^{-1}_2$ fix the subalgebra generated by $\{1,\eta\}$, we also have morphisms between $\mA^S$ and $\mB^S$. Although each of $\mA$ and $\mB$ is abelian, they do not commute with each other:
for $j\in S$, $X_j$ anticommutes with $Z_{j-1}Z_j$ and $Z_jZ_{j+1}$. We therefore define
$\mC$ to be the algebra generated by $\mA\cup\mB$ in $B(\mathcal H)$:
\[
\mC \equiv C^*(\mA\cup\mB)
\quad\text{(norm-closed $*$-algebra generated by $\mA\cup\mB$).}
\]

% We define the following linear representation on $\mC=\mA\vee \mB$:
% \begin{eqnarray}\label{represenC}
%     &&\forall a\in \mA:\pi(a)=\iota(a),\\
%     &&\forall b\in \mB:\pi(b)=\iota^{-1}(b)\ .
% \end{eqnarray}
% The Hamiltonian of the KW model is the following operator in $\mC$:
% \begin{eqnarray}
%     &&H=H_{bulk}+H_{\partial}\\
%     &&H_{bulk}=\sum_{j\in S}(X_j+Z_jZ_{j+1}),\qquad H_{\partial}=(X_L+Z_1Z_L)\ .
% \end{eqnarray}
% We refer to the first term as ``local" bulk term, whereas the second term is a ``non-local" boundary term from the point of view of operators $X_j$ with $j\in S$. The bulk Hamitlonian is manifestly invariant under $\pi$, however, the boundary term is multiplied by $\eta$:
% \begin{eqnarray}
%     &&\pi(H_{bulk})=H_{bulk}\\
%     &&\pi(H_{\partial})=\eta H_{\partial}\ .
% \end{eqnarray}
% In our example, $\eta$ is a self-adjoint unitary that generates a $\mathbb{Z}_2$ group. Since $\eta^2=1$, the operators $P_\pm=(1\pm\eta)/2$ are projections that add up to one: $P_++P_-=1$ and $P_\pm(1-P_\pm)=0$. Furthermore, since all operators in $\mC$ commute with $\eta$ the algebra decomposes as two subalgebras $\mC_\pm$ corresponding to the
% \begin{eqnarray}
%     \mC=\mC_+\vee \mC_-,\qquad \mC_\pm\equiv P_\pm\mC P_\pm\ .
% \end{eqnarray}
% Therefore, even though conjugation by $\eta$ leaves every operator invariant, $\eta$ has a non-trivial right action (equivalent to its left action) because
% any operator $c\in \mC$ has a unique decomposition $c=c_++c_-$, and $(c_++c_-)\eta=(c_+-c_-)$. The subalgebras $\mC_\pm$ have no common elements, i.e. $\mC_+\cap\mC_-=\emptyset$, and we have $\pi(H_{bulk})=H_{bulk}$, whereas for the boundary Hamiltonian we find
% \begin{eqnarray}
%     &&\pi(H_{\pm,\partial})=\pm H_{\pm,\partial}\ .
% \end{eqnarray}
% In summary, we have two non-overlapping subalgebras $\mC_\pm$ with Hamiltonians $H_\pm$  respectively, and an automorphism $\pi$ that preserves $\mC_\pm$ separately, leaves the local Hamiltonian invariant, whereas on the non-local part it acts as $\pi(H_{\pm,\partial})=\pm H_{\pm,\partial}$.

% {\bf Automorphisms of the algebra:} We define the following automorphism $\mT:\mA\to \mA$ that we call rotation (translation mod $L$):
% \begin{eqnarray}
%     &&\forall j\in \{1,\cdots, L-1\}:\qquad \mT(X_j)=X_{j+1},\qquad \mT(X_{L-1})=(X_1\cdots X_L)\eta=X_L\\
%     &&\mT(a_1a_2)=\mT(a_1)\mT(a_2),\qquad \mT(1)=1,\qquad \mT(\eta)=\eta
% \end{eqnarray}
% from which it follows that $\mT(X_L)=X_1$. Using the isometry $\iota$ and $\mT$ we extend the action of $\mT$ to $\mB$
% \begin{eqnarray}
%     \mT(\iota(a))=\iota(\mT(a))\ .
% \end{eqnarray}
% The defines a rotation $\mT$ on $\mC$ which acts as $\mT:\mC_\pm\to \mC_\pm$ because $\mT(\eta)=\eta$. Now, using $\mT$ we can define a {\it twisted isomorphism} of $\mC$:
% \begin{eqnarray}\label{twisterrepresC}
%     &&\forall a\in \mA:\tilde{\pi}(a)=\iota(a),\\
%     &&\forall b\in \mB:\tilde{\pi}(b)=\mT\circ\iota^{-1}(b)
% \end{eqnarray}
We define the following linear representation  on $\mC=\mA\vee \mB$:
\begin{eqnarray}\label{represenC}
&&\forall a\in \mA:\ \pi(a)=\iota_1(a),\\
&&\forall b\in \mB:\ \pi(b)=\iota_2(b)\color{black}\ .
\end{eqnarray}

The Hamiltonian of the KW model is the following operator in $\mC$:
\begin{eqnarray}
&&H=H_{\text{bulk}}+H_{\partial},\\
&&H_{\text{bulk}}=\sum_{j\in S}\bigl(X_j+Z_jZ_{j+1}\bigr),\qquad
H_{\partial}=\bigl(X_L+Z_1Z_L\bigr)\ .
\end{eqnarray}
There are various ways of writing this Hamiltonian:
\begin{eqnarray}
&&H=H_\mA+H_{\mB}\\
&&H_\mA=\sum_{j=1}^{L}X_j,\qquad \\
&&H_{\mB}=\iota_1(H_\mA)
%\nonumber\\
%&&H=X_1+Z_1Z_L+\sum_{j=1}^{L-1}(Z_jZ_{j+1}+\iota_2(Z_jZ_{j+1}))
\end{eqnarray}
The symmetry transformation does the following
\begin{eqnarray}
&&H_{\mA}\to \iota_1(H_{\mA})=H_\mB\\
&&    H\to H+(1-\eta)(X_2\cdots X_L)\ .
%    &&H_\mB\to \iota_2(H_{\mB})=\iota_2\circ \iota_1(H_{\mA})=\mT H_{\mA}\mT=H_{\mA}\ .
\end{eqnarray}
Another way of writing this down is
\begin{eqnarray}
H\to H+(1-\eta) (X_2\cdots X_L)(1-\eta)=H-(1-\eta)X_1(1-\eta)\ .
\end{eqnarray}

\color{brown} fix below \color{black}

We refer to the first term as a ``local'' bulk term, whereas the second term is a ``non-local'' boundary term from the point of view of operators $X_j$ with $j\in S$. The bulk Hamiltonian is invariant under $\pi$, however, the transformation of the boundary term results in an extra :
\begin{eqnarray}
&&\pi(H_{\text{bulk}})=H_{\text{bulk}},\\
&&\pi(H_{\partial})=\eta\, H_{\partial}\ .
\end{eqnarray}
\SA{We are considering the NIS operator KW to be a superoperator - might need to be careful about how we talk about the $\mathbb{Z}_2$ symmetry below}

In our example, $\eta$ is a self-adjoint unitary that generates a $\mathbb{Z}_2$ group. Since $\eta^2=1$, the operators $P_\pm=(1\pm\eta)/2$ are projections that add up to one: $P_++P_-=1$ and $P_\pm(1-P_\pm)=0$. Furthermore, since all operators in $\mC$ commute with $\eta$, the algebra decomposes via central projections as
\begin{eqnarray}
\mC=\mC_+ \oplus\ \mC_-,\qquad \mC_\pm\equiv P_\pm\,\mC\,P_\pm\ .
\end{eqnarray}
Therefore, even though conjugation by $\eta$ leaves every operator invariant, $\eta$ has a non-trivial right action (equivalent to its left action) because any operator $c\in \mC$ has a unique decomposition $c=c_++c_-$, and $(c_++c_-)\eta=(c_+-c_-)$. The subalgebras $\mC_\pm$ have trivial intersection, i.e. $\mC_+\cap\mC_-=\{0\}\color{black}$, and we have $\pi(H_{\text{bulk}})=H_{\text{bulk}}$, \color{brown} fix \color{black} whereas for the boundary Hamiltonian we find
\begin{eqnarray}
&&\pi(H_{\pm,\partial})=\pm\, H_{\pm,\partial},\qquad
\end{eqnarray}
In summary, we have two orthogonal subalgebras $\mC_\pm$ with Hamiltonians $H_\pm$ respectively, and a $*$-morphism $\pi$ that preserves $\mC_\pm$ separately, leaves the local Hamiltonian invariant, whereas on the non-local part it acts as $\pi(H_{\pm,\partial})=\pm H_{\pm,\partial}$.

{\bf Automorphisms of the algebra:} \color{brown}Add $\mathbb{Z}_2$ symmetry, full algebra of superoperators\color{black}We define the following automorphism $\mT:\mA\to \mA$ that we call rotation (translation mod $L$):
\begin{eqnarray}
&&\forall j\in\{1,\dots,L\}:\ \mT(X_j)=X_{j+1}\ \text{(indices mod $L$)},\qquad \mT(1)=1,\ \mT(\eta)=\eta
\end{eqnarray}
In particular, $\mT(X_L)=X_1$ and $\mT(\eta)=\eta$ since $\eta=\prod_{j=1}^L X_j$.
Using the morphism $\iota_1$ and $\mT$ we extend the action of $\mT$ to $\mB$ by
\begin{eqnarray}
\mT\bigl(\iota_1(a)\bigr)=\iota_1\bigl(\mT(a)\bigr)\ .
\end{eqnarray}
This defines a rotation $\mT$ on $\mC$ which acts as $\mT:\mC_\pm\to \mC_\pm$ because $\mT(\eta)=\eta$. \SA{I think the discussion of ``twisted isomorphism" only applied to the old map, here the morphisms act without twisting.}  Now, using $\mT$ we can define a {\it twisted morphism} of $\mC$:
\begin{eqnarray}\label{twisterrepresC}
&&\forall a\in \mA:\ \tilde{\pi}(a)=\iota\color{red}_1\color{black}(a),\\
&&\forall b\in \mB:\ \tilde{\pi}(b)=\sout{\mT\circ\iota^{-1}(b)}\ \color{red}\mT\circ\iota_2(b)\color{black}\ .
\end{eqnarray}
% \color{purple}
% Think about construct the operator $U_{KW}$ by enlarging $U_{\mA^S}U_{\mB^S}$, and work towards the following equations:
% \begin{eqnarray}
%     &&(U_{KW} X_j-Z_jZ_{j+1} U_{KW})P_+=0\\
%     &&(U_{KW} Z_jZ_{j+1}-X_{j+1}U_{KW})P_+=0\ .
% \end{eqnarray}
% Define the operator $D=U_{KW}P_+$, and think about the algebra
% \begin{eqnarray}
%     &&\eta^2=1, \qquad T^L=1, \qquad \eta T=T\eta,\qquad DD^\dagger=D^\dagger D=(1+\eta)\\
%     &&\eta D=D\eta=D,\qquad T^{-1}D=DT^{-1}=D^\dagger, \qquad D^2=(1+\eta)T
% \end{eqnarray}
% Note that $D$ is not self-adjoint.
% \color{black}

\SA{Similarly, in deriving $U_{KW}$ below we started from the isomorphism assumption - it seems more natural to just quote the form of the KW operator since we are not dealing with the intermediate isomorphism anymore.}


In the next section we will see that this matches the condition for \textit{recoverability} in operator algebraic QEC.

% Now we can construct the full algebra of superoperators:
%  \begin{eqnarray}
%     &&\eta^2=1, \qquad T^L=1, \qquad \eta T=T\eta,\qquad DD^\dagger=D^\dagger D=(1+\eta)\\
%      &&\eta D=D\eta=D,\qquad T^{-1}D=DT^{-1}=D^\dagger, \qquad D^2=(1+\eta)T
%  \end{eqnarray}

It is important to note that the operator implementing $X_l \rightsquigarrow Z_lZ_{l+1}$ and $Z_lZ_{l+1} \rightsquigarrow X_{l+1}$ is non-invertible specifically for this choice of algebra.  There are two ways we can avoid this:

\begin{enumerate}
\item We can work in the symmetry-restricted $\mathbb{Z}_2$ even subspace.  We can identify the $\mathbb{Z}_2$ operator with identity so that $\eta = \frac{1+\eta}{2} = 1$.  With this choice, we no longer run into issues with the kernel of the projection operator or the factor of $\eta$ in the isomorphism when acting on the boundary terms.  The KW operator can be realized unitarily as $U_{KW}$.  
\item We can work in an enlarged Hilbert space that includes defect walls corresponding to the symmetry~\cite{tachikawa2024noninvertiblysymmetryquantumoperations}.  In this Hilbert space, we can represent the action of the symmetry invertibly.
\end{enumerate}

Although the first choice is natural in many contexts where the non-invertible symmetry is implemented as a unitary composed with a projector, restricting to a symmetric subalgebra of a theory results in violations of desired algebraic properties.  We will discuss this further in continuum examples below  The second choice is a much more general principle of quantum information theory: information cannot be hidden and thus all quantum operations should be represented unitarily.  Any non-unitary operations are the result of considering some subsystem and can be realized unitarily by enlarging the system to include all relevant degrees of freedom.  The Stinespring dilation theorem makes this notion rigorous.  In the setting of this example, finite-dimensional quantum mechanics, the theorem says that any quantum operation can be represented on some enlarged Hilbert space by coupling to environmental degrees of freedom (ancilla qudits), unitary evolution of the entire system, and then tracing out the environment.  We review the Stinespring construction in Appendix \ref{sec:Stinespring}.

In this example, it is tempting to start with the property of complete positivity\footnote{$\mathrm{KW}$ is the composition of a projection and unitary, two completely positive maps, and is therefore also completely positive} and then use Stinespring dilation in the context of finite-dimensional quantum mechanics to find a unitary representation of the CP map .  The elegance of the argument in Ref.~\cite{tachikawa2024noninvertiblysymmetryquantumoperations} is in recognizing that the action of a generic (non-)invertible symmetry map on an operator is equivalent to introducing defects into the Hilbert space and moving them unitarily around the target operator $O$ realized in the defect space $\mathcal{H}_{X\overline{X}}$.  In the non-invertible case, the defects are added with an isometry $V$ that that enlarges the original Hilbert space $\mathcal{H} \rightarrow \mathcal{H}_{X\overline{X}}$ .  As such, Tachikawa's argument immediately provides a dilation $V^*\rho_{X\overline{X}}(O)V$ and uses it to prove that non-invertible symmetries are CP maps.  The prescription is completely general, has a clear physical interpretation, and can be readily applied to continuum theories.  In this work, we are interested in the non-unitary nature of the theory so we will not consider these dilations, but we recognize the representation-dependent nature of the non-unitarity.


\section{Stinespring Dilation Theorem}
\label{sec:Stinespring}

Let us review the derivation Stinespring dilation theorem and explain its significance in terms of quantum operations.  See \cite{furuya2022real, Nielsen_Chuang_2010, Heinosaari_Ziman_2011, Kretschmann_2005} for more comprehensive presentations.

First let us state the theorem in its original form:

Let $\Phi: \mathcal{A} \rightarrow B(\mathcal{H})$ be a map from a $C^*$ algebra $\mathcal{A}$ to (bounded) operators on a Hilbert space $\mathcal{H}$.  Then:
\label{eq:stinespringtheoremapp}
\begin{equation}
\Phi \;\text{completely positive}  \Leftrightarrow \Phi(a) = V^* \rho(a) V,
\end{equation}
\[  \]

where $a \in \mathcal{A}$ is a member of the $C^*$ algebra, $\mathcal{K}$ is a Hilbert space, $V: \mathcal{H} \rightarrow \mathcal{K}$ is an isometry, and $\rho: \mathcal{A} \rightarrow B(\mathcal{K})$ is a $^*$-representation.

The theorem says that we can write down the action of a CP map by moving to an enlarged Hilbert space, representing the algebra element on the new Hilbert space, and then returning to the original Hilbert space.  Note that the theorem is a two-way implication.

In the original proof, the Hilbert space $\mathcal{K}$ is constructed starting by taking an algebraic tensor product $\mathcal{A} \otimes \mathcal{H}$ of the algebra with the target Hilbert space of the map.  The resulting vector space is equipped with an inner product:
For $\alpha = \sum_ia_i\otimes x_i, \;a_i \in \mathcal{A}, \; x_i \in \mathcal{H}$ and $\beta = \sum_ib_i\otimes x_i, \;b_i \in \mathcal{A}, \; y_i \in \mathcal{H}$, \[(\alpha, \beta) \equiv \sum_{i,j}(\Phi(b_j^*a_i)x_i, y_j)\]

We quotient out $\mathcal{A} \otimes \mathcal{H}$ by $\mathcal{I}$, the set of operators $\alpha$ where $(\alpha, \alpha) = 0$.  The Hilbert space $\mathcal{K}$ is then defined as the completion of $\mathcal{A} \otimes \mathcal{H}/\mathcal{I}$ with the inner product $(\cdot\;,\;\cdot)$ defined above.

The representation $\rho$ is naturally defined by the action
\[\rho(a) \sum_jb_j\otimes x_j = \sum_j(ab_j)\otimes x_j\]

Finally, the embedding isometry $V$ acts on $\mathcal{H}$ as
\[Vx = 1 \otimes x + \mathcal{I}.\]

The Stinespring construction above is not necessarily \textit{minimal} - there could be a smaller auxillary Hilbert space such that the right hand side of Eq.~\ref{eq:stinespringtheoremapp} is still true.  However, one can show (see above references) that any two dilations are related at most by an isometry.

To understand this result physically, let us first review the notion of complete positivity in quantum information.  A map is positive if it sends positive elements to positive elements, or in the case of quantum information, valid density operators to valid density operators.  Alternatively, as is more convenient for us, we can work in the Heisenberg picture where quantum operations are applied to elements of the algebra rather than states\footnote{We can go between pictures using the dual map $\Phi_*: B(\mathcal{H}) \rightarrow \mathcal{A}$.  See Ref.~\cite{furuya2022real} Appendix A for details}.  Complete positivity means that if we enlarge the input operator by tensoring with a trivial system $a \rightarrow a \otimes \mathbbm{1}_R$ and still act with the map only on the original subsystem $(\Phi\otimes\text{id})(a \otimes \mathbbm{1}_R)$, the full system remains positive.  We impose this requirement as an axiom of quantum operations.  A typical example of a positive but not completely positive operation is the transpose map.  As an example, one can check that this map applied to the first qubit of a Bell pair does not result in a valid density matrix for the whole system~\cite{Nielsen_Chuang_2010}.

Returning to the Stinespring construction, we see that if $\Phi$ is a quantum operation from an algebra $\mathcal{A}$ to itself, (i.e. $\mathcal{A}$ is a set of operators on a Hilbert space) then the enlarged Hilbert space $\mathcal{K}$ will look like $\{\ket{a} | a\in\mathcal{A}, x \in \mathcal{H}\}$ when the null vector space is trivial.  In finite-dimensional quantum systems, the representation $\rho$ takes the simple form $\rho(a) = (a \otimes \mathbbm{1}_E)$ so that $\Phi(a) = V^*(a \otimes \mathbbm{1}_E) V$.  The subscript $E$ denotes the environmental Hilbert space which corresponds to the minimal dilation mentioned above.  We can further introduce a basis for the target space $\mathcal{H}$ and use the expression for the expectation value of operator $a$ to show that Eq.~\ref{eq:stinespringtheoremapp} takes the form \[\Phi(a) = \textrm{tr}_E(U^\dagger(a\otimes \mathbbm{1}_E)U) \] (see Section 4.2.2 of \cite{Heinosaari_Ziman_2011}).  Importantly, $U$ is a unitary on the combined space $\mathcal{H} \otimes \mathcal{H}_E$.  In practice, one usually computes the operators $U$ and $U^\dagger$ via the Kraus decomposition, which can in turn be found using the Choi map.

\section{Conditional Expectations from Fusion Categories}
\label{sec:condexp}

We start by providing a minimal review of fusion categories.  Further details can be found in the physically oriented Appendix E of Ref.~\cite{Kitaev_2006} or the mathematically oriented book Ref.~\cite{etingof2015tensor}.

A category can be thought of as a collection of objects along with morphisms between them satisfying associativity conditions.  In this context, we can think of the objects in the category $a\in \text{Obj}(\mathcal{C})$ as symmetry superoperators $\Phi_a$ and the morphisms as intertwiners between symmetry actions.  We can equip the category with further structure, including a tensor product operation $\otimes: \mC \times \mC \rightarrow \mC$ and corresponding associativity requirements.  The resulting mathematical structure is that of a \textit{fusion category}.

A key numerical invariant of an object $a\in \text{Obj}(\mC)$ is its \emph{quantum dimension} $d_a$.
Intuitively, $d_a$ measures the ``size'' of the object $a$ under fusion: for large $n$, the decomposition of
$a^{\otimes n}$ contains on the order of $d_a^{\,n}$ degrees of freedom.  Formally, in a rigid fusion category one defines the quantum dimension as the \emph{categorical trace}
of the identity,
\begin{eqnarray}
d_a \;:=\; \dim_{\mathcal C}(a)\;=\;\mathrm{Tr}_{\mathcal C}(\mathrm{id}_a),
\end{eqnarray}

In a unitary fusion category these numbers are positive real and satisfy
\begin{eqnarray}
d_a\, d_b \;=\; \sum_{c} N_{ab}^c\, d_c,
\end{eqnarray}
where $N_{ab}^c$ are the fusion multiplicities in the decomposition
\begin{eqnarray}
a\otimes b \;\cong\; \bigoplus_{c} N_{ab}^c\, c.
\end{eqnarray}

Next we define a \textit{conditional expectation} as a unital, completely positive map

\begin{eqnarray}
\mathcal{E}: \mA \rightarrow \mA^I
\end{eqnarray}

where $\mA^I$ is the algebra invariant under $\mathcal{E}$.  The conditional expectation satisfies the so-called \textit{bimodule property}

\begin{eqnarray}
\mathcal{E}(i_{1}ai_2) = i_1\mathcal{E}(a)i_2,
\end{eqnarray}

where $i_1, i_2 \in \mA^I$. This implies that the map acts is a projector.

\subsection{Rep($G$) Conditional Expectation}
\label{sec:condexprepg}

Let us verify that the map in Eq.~\ref{eq:E-def-RepG} acts like a projector onto the invariant subspace under the symmetry superoperators.

If we act with this map on an operator in a non-trivial conjugacy class (i.e. $[g]\neq 1$), we have

\begin{eqnarray}
&E_{\mC}(a_{[g]})&=\frac{1}{D^2}\sum_{c\in C}d_c\Phi_{W_c}(a_{[g]})\\
&&= \frac{1}{D^2}\sum_c (d_c \chi_c([g])) a_{[g]}\\
&&= \frac{a_{[g]}}{D^2}\sum_c (d_c \chi_c([g])) \\
&&=\frac{1}{D^2}\chi_\text{reg}([g])a_{[g]} = 0
\end{eqnarray}

since the regular representation satisfies $\chi_\text{reg}(g)= \delta_{g,e}$.

If we instead act on an operator in the trivial representation, we have

\begin{eqnarray}
&E_{\mC}(a_e)&=\frac{1}{D^2}\sum_{c\in C}d_c\Phi_{W_c}(a_e)\\
&&= \frac{1}{D^2}\sum_c (d_c \chi_c(e)) a_e\\
&&= \frac{1}{D^2}\sum_c (d_c)^2 a_e = a_e
\end{eqnarray}

where $\chi_c(e)$ is the dimension of the irrep $c$ and  $\sum_c (d_c)^2$ is the total quantum dimension squared $D^2$ by definition.

\subsection{$2d$ CFT Conditional Expectation}
\label{sec:condexpcft}

Here we show that Eq.~\ref{eq:E-def-cft} is a conditional expecation onto the invariant subalgebra.  Given some category $\mathcal{F}$ and an operator $O_b$ in representation $b$ of the chiral algebra, we have:

\begin{align}
E_{\mathcal C}(O_b)
&= \frac{1}{D_{\mathcal C}^2}\sum_{c\in \mathcal F} d_c\,\Phi_{\mathcal L_c}(O_b)\\
&= \frac{1}{D_{\mathcal C}^2}\sum_{c\in \mathcal F}
\frac{S_{1c}}{S_{11}}\frac{S_{cb}}{S_{1b}}\,O_b\\
&= \frac{O_b}{D_{\mathcal C}^2}\,\delta_{b,1}\,\frac{1}{S_{11}S_{1b}}\\
&= \frac{O_b}{D_{\mathcal C}^2}\,\delta_{b,1}\,\frac{1}{|S_{11}|^2}\\
&= \frac{O_b}{D_{\mathcal C}^2}\,\delta_{b,1}\,D_{\mathcal C}^2\\
&=
\begin{cases}
O_b, & \text{if $b$ is the identity representation},\\
0,   & \text{otherwise}.
\end{cases}
\end{align}

In the above we have used the definition of the quantum dimension in terms of the modular data $d_c = \frac{S_{1c}}{S_{11}}$, the unitarity and symmetry of the S-matrix, and the unitarity of the fusion category $\sum_a |S_{1a}|^2 = 1$.  This derivation applies broadly to examples modeled by unitary fusion categories. For example, the same equation could be used to define a passive error correcting code for anyon theories whose action on each other via linking is interpreted as an error.  The Rep($G$) conditional expectation can be considered a special case of the above derivation.


\section{Coulomb gas formalism and minimal models}

\subsection{Free boson on a cylinder}

Take the Euclidean cylinder $\Sigma=\mathbb{R}_\tau\times S^1_\sigma$ with $\sigma\sim\sigma+2\pi$.
Introduce complex coordinates
\begin{equation}
  w=\tau+i\sigma,\qquad \bar w=\tau-i\sigma,
\end{equation}
and the exponential map from cylinder to plane,
\begin{equation}
  z=e^{w},\qquad \bar z=e^{\bar w}.
\end{equation}
Let $\varphi(w,\bar w)$ be a real scalar field. We choose the standard CFT normalization
\begin{equation}
  S_0[\varphi]
  =\frac{1}{8\pi}\int d\tau\,d\sigma\;(\partial_\mu\varphi)(\partial^\mu\varphi)
  =\frac{1}{2\pi}\int d^2w\;\partial\varphi\,\bar\partial\varphi,
\end{equation}
where $\partial=\partial_w$ and $\bar\partial=\partial_{\bar w}$. On the plane this gives
\begin{equation}
  \langle \varphi(z,\bar z)\,\varphi(0)\rangle=-2\log|z|^2,
  \qquad
  \langle \varphi_L(z)\,\varphi_L(0)\rangle=-\log z,
\end{equation}
so that
\begin{equation}
  \partial\varphi_L(z)\,\partial\varphi_L(w)\sim \frac{1}{(z-w)^2}.
\end{equation}
The chiral and nonchiral vertex operators are
\begin{equation}
  V_\alpha(z)=:\!e^{i\alpha\varphi_L(z)}\!:,
  \qquad
  V_\alpha(z,\bar z)=:\!e^{i\alpha\varphi(z,\bar z)}\!:,
\end{equation}
and without background charge one has $h(\alpha)=\alpha^2/2$ (and similarly $\bar h$).

\subsection{Background charge (Feigin--Fuchs boson) and $c<1$}

To obtain $c\neq 1$, introduce a background charge $Q$ by coupling $\varphi$ to curvature:
\begin{equation}
  S_{\text{FF}}[\varphi]
  =\frac{1}{8\pi}\int_\Sigma \sqrt g\,g^{ab}\partial_a\varphi\,\partial_b\varphi
  +\frac{iQ}{8\pi}\int_\Sigma \sqrt g\,R\,\varphi.
\end{equation}
On the cylinder $R=0$, so the local dynamics is free, but on the sphere/plane the curvature term modifies
conformal Ward identities and the neutrality condition.

The holomorphic stress tensor is improved to
\begin{equation}
  T(z)= -\frac12:(\partial\varphi_L)^2:(z)+ iQ\,\partial^2\varphi_L(z),
\end{equation}
and one finds the central charge
\begin{equation}
  c = 1-12Q^2.
\end{equation}
With this $T$, the chiral vertex $V_\alpha(z)$ has conformal weight
\begin{equation}
  h(\alpha)=\frac12\,\alpha(\alpha-2Q).
\end{equation}

\subsection{The parameters $\alpha_\pm$ and degenerate charges $\alpha_{r,s}$}

Introduce $\alpha_\pm$ by
\begin{equation}
  Q=\alpha_+ + \alpha_-,
  \qquad
  \alpha_+\alpha_- = -1,
\end{equation}
in a normalization where the screening operators below have $h=1$.
Define the ``Kac momenta''
\begin{equation}
  \alpha_{r,s}=\frac12\Big((1-r)\alpha_+ + (1-s)\alpha_-\Big),
  \qquad r,s\in\mathbb{Z}_{>0}.
\end{equation}
The corresponding highest weights are
\begin{equation}
  h_{r,s}=h(\alpha_{r,s})
  = \frac{\big(r\alpha_+ + s\alpha_-\big)^2 - (\alpha_+ + \alpha_-)^2}{4}.
\end{equation}
These $h_{r,s}$ correspond to degenerate Virasoro representations with a null vector at level $rs$.

\subsection{Screenings: Virasoro-invariant integrals}

Define the screening fields
\begin{equation}
  S_+(z)=:\!e^{i\alpha_+\varphi_L(z)}\!:,
  \qquad
  S_-(z)=:\!e^{i\alpha_-\varphi_L(z)}\!:.
\end{equation}
By construction,
\begin{equation}
  h(\alpha_\pm)=1,
\end{equation}
so the screening charges
\begin{equation}
  Q_+ = \oint dz\, S_+(z),\qquad Q_-=\oint dz\,S_-(z)
\end{equation}
commute with the Virasoro algebra (they ``screen'' the background charge). In full correlators one uses
integrated holomorphic$\times$antiholomorphic screenings.

\subsection{Neutrality condition and definition of minimal-model correlators}

Because of the background charge, correlators on the sphere are nonzero only if total charge is neutralized.
For a correlator of chiral vertices $V_{\alpha_i}$ with $N_\pm$ screenings, neutrality reads
\begin{equation}
  \sum_i \alpha_i + N_+\alpha_+ + N_-\alpha_- = 2Q,
  \qquad N_\pm\in\mathbb{Z}_{\ge 0}.
\end{equation}
Minimal-model correlators are represented by free-field correlators with integrated screenings:
\begin{equation}
  \Big\langle \prod_i V_{\alpha_i}(z_i)\Big\rangle_{\text{minimal}}
  \;\propto\;
  \Big\langle \prod_i V_{\alpha_i}(z_i)\;
  \prod_{a=1}^{N_+}\int d^2u_a\, S_+(u_a)\;
  \prod_{b=1}^{N_-}\int d^2v_b\, S_-(v_b)
  \Big\rangle_{\text{free+}Q},
\end{equation}
with $N_\pm$ chosen to satisfy neutrality. The resulting screening integrals are the
Dotsenko--Fateev integrals.

\subsection{Choosing $(p,p')$: minimal models $M(p,p')$, fusion/selection rules, and an explicit Ising example}

\subsubsection{Rationality: choose $(p,p')$}
Minimal models correspond to the rational choice
\begin{equation}
  \alpha_+ = \sqrt{\frac{p'}{p}},\qquad
  \alpha_- = -\sqrt{\frac{p}{p'}},
  \qquad (p,p')=1,\ \ p,p'\ge 2.
\end{equation}
Then
\begin{equation}
  Q=\alpha_+ + \alpha_-=\sqrt{\frac{p'}{p}}-\sqrt{\frac{p}{p'}},
\end{equation}
and the central charge becomes
\begin{equation}
  c=1-6\frac{(p'-p)^2}{pp'}.
\end{equation}
The degenerate weights reproduce the Kac formula
\begin{equation}
  h_{r,s}
  =\frac{(rp'-sp)^2-(p'-p)^2}{4pp'},
  \qquad 1\le r\le p-1,\ \ 1\le s\le p'-1,
\end{equation}
with the identification $(r,s)\sim(p-r,p'-s)$. The unitary series is $p'=p+1$.

\subsubsection{Fusion rules from screening solvability}
For a correlator of degenerate primaries $\phi_{r_i,s_i}\leftrightarrow V_{\alpha_{r_i,s_i}}$,
neutrality demands integers $N_\pm\ge 0$ solving
\begin{equation}
  \sum_i \alpha_{r_i,s_i}+N_+\alpha_+ + N_-\alpha_- = 2Q.
\end{equation}
For three-point functions, solvability imposes arithmetic constraints on $(r_i,s_i)$ that reproduce
the BPZ fusion selection rules: allowed $r_3$ appear in steps of $2$ between $|r_1-r_2|+1$ and the
appropriate upper bound (and similarly for $s$).

\subsubsection{Explicit example: Ising model $M(3,4)$}
The 2d Ising CFT is the unitary minimal model with
\begin{equation}
  (p,p')=(3,4),
  \qquad
  c=1-6\frac{(4-3)^2}{12}=\frac12.
\end{equation}
We have
\begin{equation}
  \alpha_+=\sqrt{\frac{4}{3}}=\frac{2}{\sqrt3},\qquad
  \alpha_-=-\sqrt{\frac{3}{4}}=-\frac{\sqrt3}{2},\qquad
  Q=\alpha_+ + \alpha_-=\frac{2}{\sqrt3}-\frac{\sqrt3}{2}.
\end{equation}

\paragraph{(a) The three primaries $1,\sigma,\varepsilon$.}
The Kac labels are $1\le r\le 2$, $1\le s\le 3$, with $(r,s)\sim(3-r,4-s)$. The distinct primaries are
\begin{equation}
  (1,1)\sim(2,3),\qquad (1,2)\sim(2,2),\qquad (1,3)\sim(2,1).
\end{equation}
Their conformal weights are
\begin{align}
  h_{1,1}&=\frac{(1\cdot 4-1\cdot 3)^2-(1)^2}{4\cdot 3\cdot 4}=0,\\
  h_{1,2}&=\frac{(4-6)^2-1}{48}=\frac{1}{16},\\
  h_{1,3}&=\frac{(4-9)^2-1}{48}=\frac12,
\end{align}
so $\sigma\leftrightarrow(1,2)$ and $\varepsilon\leftrightarrow(1,3)$.

The corresponding Coulomb-gas charges are
\begin{equation}
  \alpha_{r,s}=\frac12\big((1-r)\alpha_+ + (1-s)\alpha_-\big),
\end{equation}
so
\begin{equation}
  \alpha_{1,1}=0,\qquad
  \alpha_{1,2}=-\frac{\alpha_-}{2}=\frac{\sqrt3}{4},\qquad
  \alpha_{1,3}=-\alpha_-=\frac{\sqrt3}{2}.
\end{equation}

\paragraph{(b) A neutrality check: $\langle \sigma\sigma\varepsilon\rangle\neq 0$.}
Consider the chiral charges
\begin{equation}
  \alpha_\sigma=\alpha_{1,2}=\frac{\sqrt3}{4},\qquad
  \alpha_\varepsilon=\alpha_{1,3}=\frac{\sqrt3}{2}.
\end{equation}
Their sum is
\begin{equation}
  2\alpha_\sigma+\alpha_\varepsilon
  =2\cdot\frac{\sqrt3}{4}+\frac{\sqrt3}{2}
  =\sqrt3.
\end{equation}
Neutrality requires
\begin{equation}
  2\alpha_\sigma+\alpha_\varepsilon + N_+\alpha_+ + N_-\alpha_- = 2Q.
\end{equation}
Compute
\begin{equation}
  2Q=2\left(\frac{2}{\sqrt3}-\frac{\sqrt3}{2}\right)=\frac{1}{\sqrt3}.
\end{equation}
Thus
\begin{equation}
  \sqrt3 + N_+\frac{2}{\sqrt3} + N_-\left(-\frac{\sqrt3}{2}\right)=\frac{1}{\sqrt3}.
\end{equation}
Multiplying by $\sqrt3$ gives
\begin{equation}
  3 + 2N_+ - \frac{3}{2}N_- = 1
  \quad\Longleftrightarrow\quad
  4N_+ - 3N_- = -4.
\end{equation}
One nonnegative integer solution is
\begin{equation}
  N_+=2,\qquad N_-=4,
\end{equation}
confirming that the Coulomb-gas construction allows a nonzero correlator in this channel (consistent with
$\sigma\times\sigma\sim \mathbf{1}+\varepsilon$).

\paragraph{(c) $\langle\sigma\sigma\sigma\sigma\rangle$ and BPZ blocks.}
The field $\sigma=\phi_{1,2}$ is degenerate at level $2$, so correlators involving $\sigma$ satisfy a
second-order BPZ differential equation. In Coulomb gas, $\langle\sigma\sigma\sigma\sigma\rangle$ is
represented as a Dotsenko--Fateev screening integral; evaluating it yields the two hypergeometric
conformal blocks corresponding to the fusion channels $\mathbf{1}$ and $\varepsilon$.

% \section{KW Operator Calculations}
% \label{sec:KWcalc}

% We would like to arrive at a unitary operator that implements this twisted isomorphism, which becomes an automorphism on all of $\mC$.  We start by considering the following operator defined on $\mC_S = \mA_S\vee \mB_S $:

% \begin{eqnarray}
% U_{\textrm{KW}^S} =\prod_{j=1}^{L-1}\lb \frac{1+i X_j}{\sqrt{2}}\rb \lb\frac{1+i Z_jZ_{j+1}}{\sqrt{2}}\rb\in \mC_S\ .
% \end{eqnarray}

% We would like to enlarge this operator to account for the boundary point at $j=L$ while respecting the rotation automorphism $\mT$.  It turns out that the correct enlargement is right-multiplying by a factor of $\frac{1+iX_L}{\sqrt{2}}$ (see Appendix A for explicit calculations).  So overall we arrive at the same operator as in Eqs.~\ref{UKWlneqL} and \ref{UKWleqL}.  Since this unitary implements the twisted isomorphisms, we have the following equations:

% For $l \neq L$,

% \begin{equation} \label{UKWlneqL}
% \begin{split}
% U_{\textrm{KW}}X_lU_{\textrm{KW}}^{\dagger} &= Z_lZ_{l+1}\\
% U_{\textrm{KW}}X_lU_{\textrm{KW}}^{\dagger} -  Z_lZ_{l+1} &= 0\\
% U_{\textrm{KW}}X_l -  Z_lZ_{l+1}U_{\textrm{KW}} &= 0\\
% (U_{\textrm{KW}}X_l -  Z_lZ_{l+1}U_{\textrm{KW}})P_+ &= 0\\
% \end{split}
% \end{equation}

% \begin{equation} \label{UKWlneqL}
% \begin{split}
% U_{\textrm{KW}}Z_lZ_{l+1}U_{\textrm{KW}}^{\dagger} - X_{l+1} &= 0 \\
% U_{\textrm{KW}}Z_lZ_{l+1} - X_{l+1}U_{\textrm{KW}} &=  0\\
% (U_{\textrm{KW}}Z_lZ_{l+1} - X_{l+1}U_{\textrm{KW}})P_+ &=  0\\
% \end{split}
% \end{equation}

% Similarly, for the boundary point,

% \begin{equation}
% \label{UKWleqL}
% \begin{split}
% U_{\textrm{KW}}X_LU_{\textrm{KW}}^{\dagger} &= \eta Z_LZ_{1}\\
% U_{\textrm{KW}}X_lU_{\textrm{KW}}^{\dagger} - \eta Z_lZ_{l+1} &= 0\\
% U_{\textrm{KW}}X_l -  Z_lZ_{l+1}U_{\textrm{KW}} \eta&= 0 \;\;\;([\eta, Z_lZ_{l+1}] = [\eta, U_{KW} ] = 0) \\
% (U_{\textrm{KW}}X_l -  Z_lZ_{l+1}U_{\textrm{KW}}\eta)P_+ &= 0  \\
% (U_{\textrm{KW}}X_l -  Z_lZ_{l+1}U_{\textrm{KW}})P_+ &= 0  \;\;\;(\eta P_+ = \eta)\\
% \end{split}
% \end{equation}
% \begin{equation}
% \label{UKWleqL}
% \begin{split}
% U_{\textrm{KW}}Z_LZ_{1}U_{\textrm{KW}}^{\dagger} &= \eta X_{1}\;\\
% U_{\textrm{KW}}Z_LZ_{1}U_{\textrm{KW}}^{\dagger} - \eta X_{1} &= 0\\
% U_{\textrm{KW}}Z_LZ_{1} -  X_{1}U_{\textrm{KW}} \eta&= 0 \;\;\;([\eta, X_1] = [\eta, U_{KW} ] = 0) \\
% (U_{\textrm{KW}}Z_LZ_{1} -  X_{1}U_{\textrm{KW}}\eta)P_+ &= 0  \\
% (U_{\textrm{KW}}Z_LZ_{1} -  X_{1}U_{\textrm{KW}})P_+ &= 0  \;\;\;(\eta P_+ = \eta)\\
% \end{split}
% \end{equation}

% The boundary and bulk operators obey this same condition.  If we define non-invertible operator  $D=U_{KW}P_+$, we find the following action:

% \begin{equation}
% \begin{split}
% DX_l&=Z_lZ_{l+1}D\\
% DZ_lZ_{l+1}&=X_{l+1}D\,,
% \end{split}
% \end{equation}

% which can be rewritten suggestively as

% \begin{equation}
% \begin{split}
% DX_l&=\tilde{\pi}(X_l)D\\
% DZ_lZ_{l+1}&=\tilde{\pi}(Z_lZ_{l+1})D\,.
% \end{split}
% \end{equation}
% \SA{This doesn't apply at the boundary point since $\iota(X_L)$ gives the extra factor of $\eta$, similar for $\iota^{-1}$}

% \color{blue}

% \underline{\textbf{Note: Why GNS does not give a unitary implementation of \text{KW}}}

% We can verify that the canonical GNS construction of the symmetric algebra generated by (including charge carrying $\textit{fields}$)

% \[\mathcal{B} = \{a\;|\;[a,\eta] = 0\} \cup \; \{b\;|\;\eta b\eta = -b\}\]

% does not lead to a unitary representation of the CP map \text{KW}.  The first set is the set $\mathcal{B}_+$ of $\mathbb{Z}_2$-even observables and the second set is the set $\mathcal{B}_{\mathcal{F}}$ of fields which connect superselection sectors of different $\mathbb{Z}_2$ charge.

% If $\mathrm{KW}$ was a *-automorphism on $\mathcal{B}$, then (Haag) section 5 theorem 2.2.4 tells us that it is implementable by a unitary $U(\mathrm{KW})$.  However, we can use the same logic as the matrix algebra case to show that it cannot be a *-automorphism:

% \underline{Proof}:

% Assume $\mathrm{KW}$ is a *-automorphism.  Then we have $\mathrm{KW}: \mathcal{B} \rightarrow \mathcal{B}$, $ \mathrm{KW}(X_j) = Z_jZ_{j+1}$, $\mathrm{KW}(Z_jZ_{j+1}) = X_{j+1}$.  Consider $\mathrm{KW}(\eta) = \mathrm{KW}(\prod_jX_j) = \prod_j\mathrm{KW(X_j)} = \prod_j(Z_jZ_{j+1}) = 1 \implies \ \eta = 1 \;\;\text{(contradiction)}$

% We can see this using the explicit GNS construction as well.  First, we choose the representative state $\rho: \mathcal{B} \rightarrow \mathbb{C}$ to be the maximally-mixed qubit state.  For this choice, there is no operator in the algebra such that $\rho(b^*b) = 0$ (since $\rho \propto \mathbf{1}$), so we have no null vector space to quotient by.  We proceed by defining an inner product as $\braket{A|B} = \omega(a^*b) = \mathrm{tr}(a^*b)$.  The GNS triple $(\pi_{\omega}, \mathcal{H}_\omega, \mathcal{H}_\omega)$ is given by the above defined Hilbert space (the original Hilbert space is already complete) and the action of the representation $\pi_\omega(a)\ket b = \ket {ab}$.  In the $C^*$-algebraic setting, $\mathrm{KW}$ is a superoperator $\mathrm{KW}: \mathcal{B} \rightarrow \mathcal{B}_+$.  In particular, we can write $\mathrm{KW}$ as a composition of a conditional expectation $\mathcal{E}_{\mathrm{KW}}: \mathcal{B} \rightarrow \mathcal{B}^+$ followed by an automorphism $\alpha_{\mathrm{KW}}: \mathcal{B}_+ \rightarrow \mathcal{B}_+$. Following Table 1, Appendix B of \cite{furuya2022real}, in the GNS Hilbert space the operator corresponding to a conditional expectation is a projection, which by definition is not unitary.

% Intuitively, this makes sense.  The GNS Hilbert space realizes a purification of some representative state $\rho$.  This purification will not in general be the same as the Stinespring dilation that realizes $\mathrm{KW}$ unitarily - this would only be the case if $\mathrm{KW}$ were a *-automorphism (rather than a *-endomorphism).

% \color{black}

\bibliography{bib}
\end{document}
